\documentclass[11pt]{article}
\usepackage[a4paper,margin=26mm]{geometry}
\usepackage{amsmath,amssymb,amsthm,hyperref}
\newtheorem{theorem}{Theorem}
\newtheorem{lemma}[theorem]{Lemma}
\newtheorem{proposition}[theorem]{Proposition}
\newtheorem{corollary}[theorem]{Corollary}
\newtheorem{remark}[theorem]{Remark}
\newcommand{\C}{\mathbb C}
\newcommand{\HH}{\mathbb H}
\newcommand{\re}{\operatorname{Re}}
\newcommand{\im}{\operatorname{Im}}
\newcommand{\Log}{\operatorname{Log}}
\title{Quadratic Kneser continuation: proofs of F and G}
\author{Research notes}
\date{4 October 2026}
\begin{document}
\maketitle
\begin{abstract}
For $f_c(w)=w^2+c$, the real critical-orbit marked Kneser crescent family
has a single-valued holomorphic continuation to the whole upper parameter
half-plane (G). Its continuation through the upper cusp corridor has the
specified sewn boundary value for every real $0<c<1/4$ (F).
The normalisation is $F_c(0)=f_c^3(0)$ throughout.
The proof combines exact crescent geometry, a marked tilted corridor,
an explicit polynomial multiplier deformation with an open-overlap
comparison, an analytic small-radius orbit estimate, and an independently
replayed Arb cover of the compact exterior core. The continuation concerns
regular height germs; a common global height domain is not asserted.
The existing two-sided Schr\"oder-based numerical iteration has no
all-upper-half-plane convergence theorem: we give an exact resonance
counterexample and state a separate, conditional local convergence result.
We then isolate a marked-cylinder continuation criterion, prove its
coordinate invariants and a general polynomial resonance obstruction,
and identify the additional hypotheses needed beyond the quadratic family.
An analytic interpolation in reciprocal degree connects the unicritical
polynomials to an exponential map. We follow this limit through the
intrinsic multiplier spectrum, local Fatou and horn coordinates, and
singular-value marked parabolic height germs.
\end{abstract}

\section{Precise targets and normalisation}
The intended exterior germ is the inverse of the two-ended, real-symmetric
crescent uniformisation, followed along the positive real critical orbit,
with $F_c(0)=w_0(c):=f_c^3(0)$ for real $c>1/4$.
The quadratic target F includes continuation through an upper cusp
neighbourhood, identification with the specified upper sewn branch for
$0<c<1/4$, and continuation throughout the upper main cardioid.
The target G is continuation of this same germ along every path in $\HH$,
jointly holomorphically near $(c,0)$ and without monodromy. A global family
with a different moving mark is not sufficient for either identification.
We distinguish a germ at $0$ from a common global time domain.

\section{Exact straight crescents on all of $\HH$}
Choose the holomorphic square root
\[
 a(c)^2=\tfrac14-c,\qquad \re a<0,\quad \im a>0,\qquad c\in\HH.
\]
Write $a=r+ib$, $b>0$, and $w=1/2+a\zeta$. The fixed points are
$L_\pm=1/2\pm a$, and the map becomes
\[
 S_a(\zeta)=\zeta+a(\zeta^2-1).
\]
The upper fixed point is $L_+$; its multiplier $1+2a$ has positive imaginary
part. The other multiplier $1-2a$ has negative imaginary part.

\begin{theorem}[Global straight-crescent geometry]\label{thm:straight}
For every $c\in\HH$, the chord $[-1,1]$ and its image $S_a([-1,1])$ meet only
at their endpoints and bound a Jordan region $H_a$. Its closed region is
parametrised by
\[
 M_a(t,s)=t+s a(t^2-1),\qquad -1\le t\le1,\quad 0\le s\le1.
\]
The open region is convex, contained in $\im\zeta<0$, and $S_a$ is injective
on the entire lower half-plane. The map $S_a$ identifies analytic collars
of the chord and its image with the correct seam orientation.
\end{theorem}
\begin{proof}
Introduce the real linear coordinates
\[
 T=\re\zeta-\frac rb\im\zeta,\qquad Y=\frac{\im\zeta}{b}.
\]
For $\zeta=M_a(t,s)$ they give $T=t$, $Y=s(t^2-1)$. Thus the open region is
exactly
\[
 -1<T<1,\qquad T^2-1<Y<0.
\]
This is an invertible linear image of the intersection of the epigraph of
the convex function $T^2-1$ with a half-plane. In particular it is convex,
and the parametrisation is bijective away from the collapsed endpoints.
Its Jacobian in the $(t,s)$ orientation is
\[
 \im(\overline{\partial_tM_a}\,\partial_sM_a)=b(t^2-1)<0.
\]
The two boundary arcs are injective and disjoint in their interiors, since
their $T$ coordinate is $t$ and their $Y$ coordinates are $0$ and $t^2-1$.
This proves the Jordan assertion.

If $S_a(\zeta)=S_a(\eta)$ and $\zeta\ne\eta$, then $\zeta+\eta=-1/a$.
The latter has positive imaginary part; hence two points in the closed
lower half-plane cannot collide. The critical point $-1/(2a)$ also has
positive imaginary part. Compactness of the chord and absence of collisions
give injective collars. Along the image arc its tangent is $1+2at$;
the determinant of this tangent with the displacement $a(t^2-1)$ is still
$b(t^2-1)<0$. The image under $S_a$ of the outer side of the chord therefore
lies on the inner side of its image. This is the required seam orientation.
\end{proof}

\begin{lemma}[Image half-plane]\label{lem:image}
In the coordinates of Theorem~\ref{thm:straight},
\[
 S_a(\{\im\zeta<0\})=\{Y<T^2-1\}.
\]
Consequently a point whose two physical preimages $w,-w$ are on the
opposite side of the chord cannot map into the lower half-plane outside
the crescent. Its first crossing to the lower half-plane is in the crescent
or on the chord.
\end{lemma}
\begin{proof}
The real line maps to the full parabola $Y=T^2-1$, with $T=t$. This curve,
completed by infinity, is a Jordan curve on the sphere. The polynomial
$S_a$ is proper, and is injective on the closed lower half-plane by
Theorem~\ref{thm:straight}. Its interior image is therefore one of the two
components of the complement of this parabola. Testing $\zeta=-i\delta$
for small $\delta>0$ selects $Y<T^2-1$.
The physical preimages of any value are $w,-w$. If neither is in the lower
half-plane, that value is outside its image. Intersecting the complement
of the image with $Y\le0$ gives the crescent, including its chord.
\end{proof}

\begin{lemma}[A uniform chord collar for the exterior core]\label{lem:collar}
Suppose $R=|a|\le1/2$, $-R^2\le r\le0$, $b>0$. With $\eta=1/4$, the region
\[
 K_a=\{\zeta=x+iy:|x|<1,\ |y|<\eta b(1-x^2)\}
\]
is a holomorphic chord chart of the quotient: its lower part is inside
$H_a$, and its upper part maps inside $H_a$ under $S_a$. The two
representations are identified by the seam.
\end{lemma}
\begin{proof}
Here $b^2\ge3R^2/4$, $|r|\le R^2\le1/4$. For $y<0$, the margin above
the image parabola in the coordinates of Theorem~\ref{thm:straight} is
\[
 Y-(T^2-1)=1-x^2+\frac yb(1+2rx)-\frac{r^2y^2}{b^2}
 >(1-x^2)(1-3\eta/2-\eta^2/16)>0.
\]
For $y>0$, write $T',Y'$ for the coordinates of $S_a(x+iy)$. Directly,
\[
 T'=x-(r+2R^2x)y/b,\quad
 Y'=x^2-y^2-1+(1+2rx)y/b,
\]
and
\[
 Y'-(T'^2-1)=\frac yb|1+2ax|^2
             -y^2\left(1+\frac{(r+2R^2x)^2}{b^2}\right)>0.
\]
Indeed the squared modulus is at least $3/4$, the parenthesis at most $4$,
and $4by\le\eta$. Also $Y'<-(1-x^2)(1-3\eta/2)<0$.
These inequalities imply $|T'|<1$ and $S_a(x+iy)\in H_a$.
Finally $S_a$ is injective on $\im\zeta<\eta b$, since the imaginary part
of its collision sum $-1/a$ is $b/R^2\ge4b>2\eta b$.
Thus this is an analytic seam collar, not an approximate membership test.
\end{proof}

\begin{theorem}[An explicit global cylinder model]\label{thm:qc}
Let $X_a$ be $H_a$ with the two open edges identified by $S_a$, and with the
fixed endpoints omitted. Then $X_a\simeq\C^*$.
The family has a single-valued, jointly holomorphic uniformisation near
every interior holomorphic marked section. In particular
\[
 m(c)=\tfrac12-\tfrac12 a(c)^2
\]
defines such a section on all of $\HH$.
\end{theorem}
\begin{proof}
Parametrise $t=\tanh x$ and use the source cylinder coordinate $z=s+ix$,
where $x\in\mathbb R$ and $s\in[0,1]$, identifying $s=0$ and $s=1$.
The map to the physical plane is $J_a=1/2+aM_a(\tanh x,s)$.
Since the affine factor $a$ cancels from the dilatation, put $B=1+2sat$.
The Beltrami coefficient is
\[
 \mu_a(x,s)=\frac{-a+iB}{-a-iB}.
\]
An exact calculation gives
\[
 |a+iB|^2-|-a+iB|^2=4b>0.
\]
Consequently $|\mu_a|<1$, uniformly for $(t,s)\in[-1,1]\times[0,1]$
at each fixed parameter, and uniformly on compact parameter sets. The
coefficient is holomorphic in $a$ as an $L^\infty$-valued function.
Extend it periodically in $s$. The edge identification agrees with the
holomorphic seam map, so this measurable coefficient gives the pulled-back
complex structure across the seam as well.

The map is a quasiconformal homeomorphism from $\C/\mathbb Z$ to $X_a$.
Under $\xi=e^{2\pi iz}$ its ends are $0,\infty$; quasiconformal
uniformisation of the sphere, after assigning arbitrary values to the
coefficient at these two points, proves $X_a\simeq\C^*$.
The $x\to+\infty$ end corresponds to $L_+$ and hence to $\xi=0$.

For the last assertion apply the analytic parameter version of the
Ahlfors--Bers theorem to this coefficient. Initially normalise its solution
$W_a$ at $0,1,\infty$. In a fixed smooth interior source chart, both
$V_a=W_a\circ\exp(2\pi iz)$ and $J_a$ are holomorphic in $a$ at fixed source
point and have the same Beltrami coefficient. Put $x=J_a^{-1}(w)$, using
$x$ here for a two-dimensional source coordinate. Differentiating
$J_a(x)=w$ at fixed $w$ gives
$x_{\bar a}+\mu_a\bar x_{\bar a}=0$; hence
\[
 \partial_{\bar a}(V_a\circ J_a^{-1})
 =(V_a)_x(x_{\bar a}+\mu_a\bar x_{\bar a})=0.
\]
The solutions are $C^1$ in these smooth interior charts by interior elliptic
regularity. Thus the descended coordinate is jointly holomorphic in the
physical variable and parameter there. In a physical seam collar it is
continuous in both variables by continuous dependence of the normalised
Beltrami solutions. Off the analytic seam hypersurface it is jointly
holomorphic, either by this interior computation or after the holomorphic
chart change $w\mapsto f_c(w)$. The continuous values coincide on the
seam by the edge identification. Morera's theorem, applied to transverse
complex lines and then by continuity to the remaining lines, makes that
analytic seam removable. Joint holomorphy therefore holds in the seam
collars too. Multiplication by its nonzero value at
an interior holomorphic section normalises that section to $1$.
Uniqueness with the two ordered ends and this mark patches all local charts.
Finally $m(c)=1/2+aM_a(0,1/2)$ is interior for every $c$.
\end{proof}

\begin{remark}[What the global model does not prove]
Normalising at $m(c)$ gives a global family of inverse-uniformisation germs
with that value at time $0$. It does not identify the critical-orbit marked
exterior family. A change to a Fourier constant mode likewise requires a
proved conversion back to the specified mark. This distinction is essential.
\end{remark}

\section{The outer region with the original normalisation}
\begin{theorem}[The outer parameter region]\label{thm:outer}
The real critical-orbit marked crescent family continues to a single-valued
holomorphic family of germs with $F_c(0)=f_c^3(0)$ on
\[
 \mathcal O=\{c\in\HH:|c-1/4|>1/4\}.
\]
It agrees with the exterior real family near every real $c>1/2$.
\end{theorem}
\begin{proof}
Mark the quotient at the physical critical value $w=c$. Set $R^2=|a|^2$.
At this value the coordinates of Theorem~\ref{thm:straight} are
\[
 T(c)=-\frac{r}{2R^2},\qquad Y(c)=-1+\frac1{4R^2}.
\]
If $R>1/2$, then $|T(c)|<1$, $Y(c)<0$, and
\[
 Y(c)-(T(c)^2-1)=\frac{b^2}{4R^4}>0.
\]
Thus $c$ is strictly inside the crescent on $\mathcal O$. Use
Theorem~\ref{thm:qc} to normalise the quotient coordinate at this mark, and
let $U_c$ be its locally inverse logarithmic lift, $U_c(0)=c$.
Crossing the seam gives $U_c(z+1)=f_c(U_c(z))$ in the successive dynamical
charts. Define the requested germ by
\[
 F_c(z)=f_c^2(U_c(z)).
\]
This is jointly holomorphic and has $F_c(0)=f_c^3(0)$; it obeys the
functional equation on its continued charts. All choices use a single
global marked quotient, so there is no parameter monodromy on $\mathcal O$.

For real $c>1/2$, the physical crescent intersects the real axis in
$(1/2,c+1/4)$, which contains $c$. Reflection interchanges the two ends,
fixes this mark, and makes the logarithmic coordinate real along the
real interval. The seam gives increase $1$ under $f_c$. The resulting
inverse is precisely the real symmetric crescent construction, followed
forward twice to the specified value $f_c^3(0)$.
\end{proof}

\begin{proposition}[Local parameter analyticity on the whole exterior ray]
The same real family has a jointly holomorphic parameter germ near every
real $c>1/4$, with the normalisation $F_c(0)=f_c^3(0)$.
\end{proposition}
\begin{proof}
For real $c>1/4$, $f_c(x)-x=(x-1/2)^2+c-1/4>0$. Its critical orbit
$w_n=f_c^n(0)$ is strictly increasing and unbounded, since a finite limit
would be a real fixed point. The first crossing of $1/2$ lands in
$[1/2,c+1/4]$. At a strict crossing its orbit point is interior to the
real crescent; at equality use the analytic seam collar. Hence some finite
orbit chart is a mark, and the same chart persists in a complex parameter
neighbourhood. Pulling back from this chart to $w_3$, or pushing forward
if its index is below $3$, gives the required germ. Every inverse step is
at a strictly positive orbit point, where $f_c'$ is nonzero. Reflection and
the ordered ends give the same real germ; uniqueness then patches the
local parameter discs along the ray.
\end{proof}

\section{Critical-orbit charts on the complete exterior}
Put $w_n=f_c^n(0)$ and
\[
 P_a(w)=b\re w-r\im w.
\]
The chord is $P_a(w)=b/2$; its crescent side is $P_a(w)>b/2$.
Both $w$ and $-w$ are on the opposite side if $0<P_a(w)<b/2$.

\begin{lemma}[Analytic small-radius entry]\label{lem:tail}
Let $0<R\le1/16$, $0\le s\le1$, and
\[
 a=-R^2s+iR\sqrt{1-R^2s^2},\qquad c=1/4-a^2.
\]
The critical orbit first crosses the chord in at most $\lceil2/R\rceil$
steps. That crossing is in $H_a$ or on its chord. Every earlier $w_j$,
$j\ge1$, has positive real part and $0<P_a(w_j)<b/2$.
\end{lemma}
\begin{proof}
Write $w_n-1/2=x_n+iy_n$, initially $x_0=-1/2$, $y_0=0$. Exactly,
\[
 x_{n+1}=x_n+x_n^2+R^2-2R^4s^2-y_n^2,\qquad
 y_{n+1}=(1+2x_n)y_n+2R^3s\sqrt{1-R^2s^2}.
\]
As long as $-1/2\le x_n<0$ and $n\le\lceil2/R\rceil$, induction gives
$0\le y_n\le2R^3n\le5R^2$. Consequently
\[
 x_{n+1}-x_n\ge x_n^2+(1-27R^2)R^2
 \ge x_n^2+\tfrac{89}{100}R^2>0.
\]
If also $x_{n+1}<0$, integration of the derivative of $\arctan(x/R)$
gives an increase at least $89R/100$: its denominator $R^2+t^2$ is
at most $R^2+x_n^2$ along this negative interval. Remaining negative for
$\lceil2/R\rceil$ steps would give an increase at least $1.78>\pi/2$,
a contradiction. This also closes the induction used in the bounds.
Thus some $x_N\ge0$ occurs by this time. Before then $\re w_n>0$ for
$n\ge1$, $y_n\ge0$, and
$P_a(w_n)=b(x_n+1/2)+R^2sy_n>0$.
At $x_N\ge0$ this quantity is at least $b/2$, so a first chord crossing
occurs no later. At its preceding point both physical preimages are on
the opposite side. Lemma~\ref{lem:image} places its image in the crescent
or on the chord; all preceding inverse steps are noncritical.
\end{proof}

\begin{lemma}[Certified compact core]\label{lem:core}
For the same parametrisation, all $(R,s)\in[1/16,1/2]\times[0,1]$ have
an orbit chart of index $1\le n\le21$ in $H_a$ or $K_a$. For every
$1\le j<n$, $\re w_j>0$ and $0<P_a(w_j)<b/2$.
\end{lemma}
\begin{proof}
This is a finite interval-arithmetic assertion. The accompanying
\texttt{certify\_quad\_entry.py} starts with the three closed rectangles
$[1/4,1/2]$, $[1/8,1/4]$, $[1/16,1/8]$ times $[0,1]$.
It uses rational bisections and Arb outward-rounded balls at 60 decimal
digits. On each rectangle it propagates $w\mapsto w^2+c$ and verifies
either the strict inequalities of Theorem~\ref{thm:straight} or those of
Lemma~\ref{lem:collar}; at every earlier index it verifies the two
displayed strict inequalities for $P_a$ and $\re w>0$.
All 4,013 final rectangles passed: 3,092 use the interior and 921 the
collar. There are no pending rectangles. A separate program,
\texttt{verify\_quad\_entry.py}, imports no generator code: it replays
every asserted orbit and predicate at 90 decimal digits and checks the
rectangle cover by an exact rational sweep. On each of 282 consecutive
$R$ slabs the active closed $s$ intervals tile $[0,1]$ once. Their
endpoints cover all slab boundaries, proving coverage of the closed
rectangle, rather than a grid of sample points.
The replay result is \texttt{quadratic-entry-verification.json}; it
records the SHA-256 of \texttt{quadratic-entry-core.json} and maximum
index 21. These finite checks and the analytic collar lemma prove the
assertion.
\end{proof}

\begin{lemma}[Compatibility of first-entry charts]\label{lem:entry-patch}
The orbit charts in Lemmas~\ref{lem:tail} and~\ref{lem:core}, and the
critical-value chart of Theorem~\ref{thm:outer}, give compatible germs
with value $w_3$ at time zero.
\end{lemma}
\begin{proof}
An interior mark $w_n$ has $P_a(w_n)>b/2$. A collar mark either has this
property or is just above the chord, in which case $f_c(w_n)$ is interior
by Lemma~\ref{lem:collar}. Thus, if two allowed charts are used at one
parameter, their indices differ by at most one: a later index would
violate the required preceding inequalities. In the one-step case the
two physical marks are related by the analytic seam $w\sim f_c(w)$.
The same statement holds in a small open parameter neighbourhood,
using the interior charts and analytic seam collars. In particular the
core collar inequalities are only used on their certified domain. At an
upper collar mark $S_a(\zeta)$ is strictly interior and remains so on
a parameter disc by continuity; at the chord itself one uses the analytic
seam. We do not apply Lemma~\ref{lem:collar} to parameters outside its
stated domain in order to cross the neutral arc.

Normalise the logarithmic lift at a chart of index $n$ to have time $n-3$.
Its inverse $U$ is transported back to $w_3$ by the selected inverses
of $f_c$, or forward when $n<3$. All required inverse steps have
$\re w_j>0$, so their local branches are noncritical and agree on
overlaps. An index change by one adds one seam step and removes exactly
one inverse step. The ordered ends and the mark determine the quotient
uniformisation uniquely; hence the resulting germs agree on these
open overlaps, including at a chord crossing.

At $R=1/2$, the critical value is on the chord, with
$T(c)=-r/(2R^2)$. On the exterior part this lies in $[0,1/2]$.
Its collar therefore bridges the interior-mark construction for $R>1/2$
and the core construction for $R\le1/2$. On the overlap it is the same
seam marking of $w_1$, followed forward twice. This proves compatibility
also with the outer family. Subsequent physical returns of an orbit to
the crescent are never used in this argument.
\end{proof}

\begin{theorem}[The complete exterior and every open neutral arc point]
\label{thm:exterior}
Let $\mathcal C=\{\lambda/2-\lambda^2/4:|\lambda|<1\}$ be the main
cardioid. The original exterior family extends to a single-valued
holomorphic family of germs on $\HH\setminus\overline{\mathcal C}$,
and across every point of $\partial\mathcal C\cap\HH$.
\end{theorem}
\begin{proof}
The upper multiplier is $\lambda=1+2a$, and
$|\lambda|^2=1+4r+4R^2$. For $R\le1/2$, being exterior or neutral is
equivalent to $-R^2\le r<0$, precisely the parametrisation of
Lemmas~\ref{lem:tail} and~\ref{lem:core}, with $s=-r/R^2$.
They supply all its marks; Theorem~\ref{thm:outer} supplies every $R>1/2$.
At a tail crossing on the chord use its analytic collar. All these marks
persist on open parameter discs, and Theorem~\ref{thm:qc} and
Lemma~\ref{lem:entry-patch} give the holomorphic family and its gluing.

For a neutral point write $\lambda=e^{i\theta}$, $0<\theta<\pi$.
Then $R=\sin(\theta/2)$ and $r=-R^2$. For $R\le1/16$ the analytic
lemma applies, for $1/16\le R\le1/2$ the certificate applies, and
for $R>1/2$ the critical value is strictly interior. The open chart
discs cross the arc, including every rational and irrational neutral
multiplier. End linearisation was not used.
Agreement with the real exterior ray follows from its local first-entry
construction, proved above, and the same patching rule.
\end{proof}

\section{Explicit cusp estimates}
Put $c=1/4-\varepsilon^2/4$ and $w=1/2+u$. Then
\[
 P_\varepsilon(u)=u+u^2-\varepsilon^2/4,\quad
 u_1=-\varepsilon/2,\quad u_2=\varepsilon/2,\quad
 \lambda_1=1-\varepsilon,\quad\lambda_2=1+\varepsilon.
\]
The map $\varepsilon\mapsto c$ is injective on the sector
$-7\pi/12<\arg\varepsilon<\pi/12$. Its image contains an upper punctured
half-disc at $1/4$ and both adjacent real segments.
Set
\[
 q=\frac{u+\varepsilon/2}{u-\varepsilon/2},\quad \chi=\log q,
 \quad u=\frac{\varepsilon}{2}\frac{e^\chi+1}{e^\chi-1},
 \quad \mathcal G=\{\operatorname{dist}(\chi,2\pi i\mathbb Z)\ge1/2\}.
\]

\begin{lemma}[Step estimates with explicit constants]\label{lem:step}
For $0<\rho=|\varepsilon|\le1/128$ on $\mathcal G$,
\[
 |u|\le6\rho,\qquad
 \Delta(\chi)=\Log(1+u-\varepsilon/2)-\Log(1+u+\varepsilon/2),
\]
represents the step $\chi(P_\varepsilon u)=\chi(u)+\Delta$ and satisfies
\[
 |\Delta+\varepsilon|\le7\rho^2,\qquad
 |\partial_\chi\Delta|\le40\rho^2.
\]
If $\omega=e^{i\alpha}$, $\cos\alpha\ge3/5$, and
$-\im(\varepsilon\bar\omega)\ge\rho/2$, then
\[
 \tfrac25\rho\le\im(\Delta\bar\omega)\le\tfrac65\rho.
\]
\end{lemma}
\begin{proof}
On $\mathcal G$, $|1-e^\chi|\ge\max(1,|e^\chi|)/6$.
For completeness reduce modulo $2\pi i$ to $|\im\chi|\le\pi$.
The function $(e^\chi-1)^{-1}$ has modulus at most $1$ on the horizontal
edges, limits $-1,0$ at the two infinite ends, and modulus at most
$(1-e^{-1/2})^{-1}$ on $|\chi|=1/2$. Indeed, if
$v=e^\chi-1$ has $|v|<1$, the principal logarithm gives
$1/2=|\Log(1+v)|\le-\log(1-|v|)$; if $|v|\ge1$ the bound is immediate.
The maximum principle on the exterior of this disc gives
$|1-e^\chi|\ge1-e^{-1/2}$; for $\re\chi\ge\log2$ the sharper bound
$|1-e^\chi|\ge|e^\chi|/2$ applies. Together with reflection these imply the
stated bound with $1/6$. Thus $|u|\le6\rho$ and
$|u_\chi|=|\varepsilon e^\chi/(e^\chi-1)^2|\le36\rho$.
The exact factorisation of $P_\varepsilon(u)-u_j$ gives the formula for
the ratio of the two cross-ratios. Both logarithms here are the branches
near $1$, since $|u\pm\varepsilon/2|\le13\rho/2<1$.
Also
\[
 \Delta=-\varepsilon\int_{-1/2}^{1/2}\frac{dt}{1+u+t\varepsilon}.
\]
It follows that
\[
 |\Delta+\varepsilon|\le
 \frac{13/2}{1-13\rho/2}\rho^2<7\rho^2,\quad
 |\Delta_\chi|\le\frac{36}{(1-13\rho/2)^2}\rho^2<40\rho^2.
\]
The normal-component bounds are $\rho/2-7\rho^2$ and $\rho+7\rho^2$,
which imply the claimed rational inequalities.
\end{proof}

\begin{proposition}[Tilted initial regions and their flat quotient]
For every line $\ell=i\pi+x_0+\mathbb R\omega$ as in
Lemma~\ref{lem:step}, $|x_0|\le1/2$, the line and its image under
$\chi\mapsto\chi+\Delta$ bound a thin strip. Its parametrisation is
\[
 N(t,s)=i\pi+x_0+t\omega+s\Delta(i\pi+x_0+t\omega).
\]
It is embedded by $\chi\mapsto u(\chi)$ as a Jordan crescent. The edge
quotient is $\C^*$, including at neutral multipliers, and admits a flat
model with Beltrami bound
\[
 \|\mu\|_\infty\le\frac{25\rho}{1-25\rho}<1.
\]
\end{proposition}
\begin{proof}
The distance from $\ell$ to every hole is at least
$\pi\cos\alpha-|x_0\sin\alpha|\ge3\pi/5-2/5>1.48$.
Thus its $1/4$-strip is in $\mathcal G$.
The derivative estimate in Lemma~\ref{lem:step} makes
$\chi+\Delta$ injective on this convex strip and gives a positive real
derivative for the tangential coordinate of the image of $\ell$.
The image is therefore a graph; its normal height is in
$[2\rho/5,6\rho/5]$. The Jacobian of $N$ is positive, since it is at least
$2\rho/5-40\rho^2(\rho+7\rho^2)>0$.
The parametrisation is proper and bijective on both boundary lines;
it is therefore a diffeomorphism onto the strip. Its width is less than
$1/4$. A $2\pi i$-translate changes the normal coordinate by at least
$6\pi/5$, so the exponential, and hence $u(\chi)$, is injective there.
The two infinite ends tend to $u_1,u_2$, establishing the Jordan assertion.

Put $A=1/\Log(1-\varepsilon)$ and compare $AN(t,s)$ with the straight strip
$A(i\pi+x_0+t\omega)+s$. The integral estimates above and
$|\Log(1-\varepsilon)+\varepsilon|\le\rho^2/[2(1-\rho)]$ give
$|A\Delta-1|\le8\rho$. In the unit tangential direction the derivative
error is at most $40\rho^2$; in the $s$ direction it is at most $8\rho$.
The sine of the angle between these two source directions is at least
$49/100$. The inverse basis matrix has norm less than $3$, so the
derivative error has norm at most $3(8\rho+40\rho^2)\le25\rho$.
The displayed Beltrami bound follows from the Wirtinger derivatives.
The map respects the edge identifications, and the resulting
quasiconformal cylinder again has two punctured ends. This proves the
quotient assertion without requiring linearisation at a neutral endpoint.
\end{proof}

\section{Polynomial multiplier transport and its comparison}
\begin{lemma}[An explicit holomorphic multiplier deformation]\label{lem:surgery}
Let $\lambda_*=i/2$, $c_*=1/16+i/4$, and $p_*=i/4$.
There is a single-valued holomorphic motion $h_\lambda$ on
\[
 D_\lambda=\{0<|\lambda|<1:-\pi<\arg\lambda<\pi\}
\]
with $h_{\lambda_*}=\mathrm{id}$, $h_\lambda(0)=0$, and
\[
 h_\lambda\circ f_{c_*}=f_{c(\lambda)}\circ h_\lambda,
 \qquad c(\lambda)=\lambda/2-\lambda^2/4.
\]
It transports the marked crescent at $c_*$ to a two-ended cylinder.
Its inverse-uniformisation, followed forward twice, defines a
single-valued jointly holomorphic family $F^{\rm tr}$ near time zero,
with $F^{\rm tr}_{c(\lambda)}(0)=f_{c(\lambda)}^3(0)$.
\end{lemma}
\begin{proof}
The upper attracting fixed point is $p_*$ and
$c_*-p_*=1/16$. On $|v|\le1/16$,
\[
 f_{c_*}(p_*+v)-p_*=(i/2)v+v^2,\qquad
 |(i/2)v+v^2|\le(9/16)|v|.
\]
A nonzero $v$ in this disc never maps to zero. The disc avoids the
critical point. Thus the critical value is in the attracting basin
$B_*$, and its orbit never hits $p_*$ or the critical point.
Extend the Koenigs function $\sigma$ to the whole basin by
$\sigma(w)=\lambda_*^{-n}\sigma(f_{c_*}^n(w))$ once the orbit reaches
the linearisation disc. This is well-defined and holomorphic, with
$\sigma\circ f_{c_*}=\lambda_*\sigma$ and
$\sigma(c_*),\sigma'(c_*)\ne0$. Its zero and derivative-zero sets
are discrete in each basin component. Also $\sigma(-w)=\sigma(w)$.
The full basin is bounded and completely invariant.

Put $s=\Log\lambda$, $s_*=\Log\lambda_*$, and
\[
 \beta=\frac{s-s_*}{2\re s_*},\qquad
 k=\frac{s-s_*}{s+\overline{s_*}},\qquad
 \phi_\lambda(\zeta)=\zeta|\zeta|^{2\beta}.
\]
Since $\re s,\re s_*<0$, $|k|<1$, uniformly on compact parameter sets.
The radial exponent has positive real part, so $\phi_\lambda$ is an
orientation-preserving homeomorphism, and
$\phi_\lambda(\lambda_*\zeta)=\lambda\phi_\lambda(\zeta)$.
Define, almost everywhere,
\[
 \nu_\lambda(w)=
 k\frac{\sigma(w)}{\overline{\sigma(w)}}
   \frac{\overline{\sigma'(w)}}{\sigma'(w)}\quad(w\in B_*),
 \qquad \nu_\lambda=0\quad(w\notin B_*).
\]
Assign zero on the discrete exceptional sets. Differentiating the
Koenigs identity shows
$f_{c_*}^*\nu_\lambda=\nu_\lambda$ almost everywhere. This includes
the precritical set: no division at a critical point is asserted.
The coefficient is even, has bounded support, and is holomorphic in
$\lambda$ as an $L^\infty$ function.

Let $h_\lambda$ be its principal solution, normalised by
$h_\lambda(w)=w+O(1/w)$ at infinity. Ahlfors--Bers gives holomorphic
dependence at every fixed point. Uniqueness applied to
$-h_\lambda(-w)$ makes it odd, hence $h_\lambda(0)=0$.
Invariance makes $h_\lambda f_{c_*}h_\lambda^{-1}$ holomorphic.
It is a proper degree-two map of the plane, monic at infinity and
critical at zero, so it is $w^2+\widetilde c$.
Near $p_*$, $\phi_\lambda\circ\sigma\circ h_\lambda^{-1}$ is a
conformal linearising coordinate, giving multiplier $\lambda$.
Its fixed point must be $\lambda/2$, so
$\widetilde c=\lambda/2-\lambda^2/4$. The conjugacy therefore has
exactly the asserted polynomial normalisation.

Here $|a(c_*)|^2=5/16>1/4$, so $c_*$ is an interior mark of the
straight crescent. Transport that crescent and its seams by $h_\lambda$.
The pulled-back structure on its fixed quotient is the restriction
of $\nu_\lambda$, well-defined by invariance and of norm $<1$.
Uniformising this fixed quotient with its two ordered ends and the
critical-value mark produces its transported $\C^*$ family.
Near $c_*$ the coefficient is smooth, because $\sigma,\sigma'$ are
nonzero there. Both the physical motion and the quotient Beltrami
solution are holomorphic in the parameter at fixed source point and
solve the same equation. The cancellation computation in
Theorem~\ref{thm:qc} therefore gives joint holomorphy in the physical
variable and parameter near $h_\lambda(c_*)=c(\lambda)$.
The holomorphic inverse function theorem and a forward push by $f_c^2$
give $F^{\rm tr}$ with the asserted mark. No pullback through zero
is involved.
\end{proof}

\begin{lemma}[Open-overlap comparison]\label{lem:transport-compare}
On an open neighbourhood of $c_*=1/16+i/4$, the transported family
equals the original marked straight-crescent family.
\end{lemma}
\begin{proof}
Use the fixed-point-pair logarithm
\[
 \chi=\log\frac{w-L_+}{w-L_-}.
\]
The chord lifts to $\ell=\pi i+\mathbb R$, and its image to
$\chi\mapsto\chi+\Delta_c(\chi)$. At the base parameter
$a_*=-1/2+i/4$, for a chord point $\zeta=t\in[-1,1]$,
\[
 \Delta_{c_*}=\Log(1+a_*(t+1))-\Log(1+a_*(t-1)),\qquad
 \frac{d\chi(f(w))}{d\chi(w)}
 =\frac{1+2a_*t}{(1+a_*t)^2-a_*^2}.
\]
The imaginary part of the step is strictly positive and uniformly
bounded: the numerator in the argument of its ratio has imaginary
part $1/2$ and real part $11/16-t+5t^2/16\ge0$.
Moreover the real part of the derivative is
\[
 \frac{16(-5t^3+23t^2-27t+13)}
 {(5t^2-26t+37)(5t^2-6t+5)}>0.
\]
The cubic numerator has Bernstein coefficients
$68,28/3,4/3,4$ on $t\in[-1,1]$, all positive; both denominator
factors are positive there. Hence the image is a graph and the lifted
crescent lies in a horizontal band between $\pi$ and $\pi+M$, with
step normal component at least some $m>0$. These inequalities persist
near the base, on a small analytic collar. At the two infinite ends
the step tends uniformly to $\Log\lambda$ and $-\Log(2-\lambda)$,
whose imaginary parts remain positive. Thus, for sufficiently large
$|\re\chi|$, the same normal lower bound holds for every $\im\chi$.

We justify comparison at the infinite ends, where compact convergence
of the motion alone would be insufficient. Both fixed points at the
base are hyperbolic and noncritical. In their parameter-dependent
Koenigs coordinates the lifted conjugacy satisfies
\[
 H_\lambda(v+\Log\lambda_*)=H_\lambda(v)+\Log\lambda,
 \qquad H_\lambda(v+2\pi i)=H_\lambda(v)+2\pi i
\]
at the attracting end, and the corresponding equations with
$2-\lambda_*$, $2-\lambda$ at the repelling end.
On a fixed compact fundamental annulus the motion and its lift tend
uniformly to the identity. Writing a point as an integer iterate of
that annulus gives, with $\delta\to0$ as $\lambda\to\lambda_*$,
\[
 |H_\lambda(v)-v|\le C\delta(1+|\re v|).
\]
Koenigs coordinates differ from the pair logarithm by constants and
exponentially small terms at the ends. Consequently the entire lifted
transported crescent lies within a tube of width
$C\delta(1+|\re\chi|)$ about the lifted reference crescent.

At a far end, repeatedly apply the lifted map or its inverse until
the point reaches the reference crescent, choosing the first crossing
of its chord. The normal component is at least $m/2$ and each step
has modulus at most $2M$. Thus at most
$C_1\delta(1+|\re\chi|)+C_2$ steps suffice. Their total displacement
has the same bound. Take $\delta$ small enough that their real parts
stay in the chosen far end. No hole or critical point is then met.
On the compact middle part, convergence of the motion and the
injective analytic seam collars give the same construction using
finitely many charts. The closed reference crescent avoids the physical
critical point $0$, so these compact collars may all be chosen
noncritical; at the far ends this follows from their fixed-point
neighbourhoods. Every chart map is a fixed noncritical iterate
or inverse iterate of $f_c$; an index change at a reference edge is
exactly its seam identification. Thus these maps descend to a
holomorphic map of the transported quotient into the reference
quotient.

This map is proper at both ends: the displacement bound is less than
half of $|\re\chi|$ sufficiently far out. It preserves their ordering
and sends one seam circuit to one seam circuit. After uniformisation
it is a proper map $\C^*\to\C^*$ extending to the compactified spheres,
with only one zero and one pole, and with induced winding number one.
Its degree is therefore one, and it is biholomorphic. The critical
value is interior to both crescents near the base and is represented
without an iterate, so the map preserves the time-zero marking before
the common forward push by two. Uniqueness of the three-point
uniformisation gives equality of the germs on an open parameter set.
\end{proof}

\begin{theorem}[Quadratic G]\label{thm:G}
The critical-orbit marked real Kneser crescent family for $c>1/4$
has a single-valued holomorphic continuation to all of $\HH$.
For each $c_0\in\HH$ it is represented by a jointly holomorphic
$F(c,z)$ near $(c_0,0)$, with $F(c,0)=f_c^3(0)$. In particular its
continuation along every upper-half-plane parameter path is independent
of the path.
\end{theorem}
\begin{proof}
The map $\lambda\mapsto c(\lambda)$ is injective on $|\lambda|<1$:
two different preimages would have sum $2$, impossible in this disc.
Also $\im c(\lambda)=\im\lambda(1-\re\lambda)/2$.
Thus the upper half of the main cardioid is the image of
$\{|\lambda|<1,\im\lambda>0\}$, and
Lemma~\ref{lem:surgery} supplies a single-valued family there.
Lemma~\ref{lem:transport-compare} identifies it with the outer family
on an open subset. The intersection of the cardioid with
$|c-1/4|>1/4$ is connected: in multiplier coordinates it is
$|1-\lambda|>1$, $|\lambda|<1$, $\im\lambda>0$.
Horizontal slices are intervals on the left of the circle centred at
$1$; these slices join continuously. The identity theorem therefore
identifies the families on this entire overlap.

Theorem~\ref{thm:exterior} gives compatible open chart discs along the
whole neutral arc. Starting where $R>1/2$, their inner portions agree
with $F^{\rm tr}$. Any compact subarc is covered by a finite chain of
overlapping discs centred on the arc; consecutive inner portions have
an open intersection. The entry-patching lemma and the identity theorem
propagate the equality along this chain. Every open arc point is reached
by such a compact chain; the cusp and the real endpoint are excluded
from $\HH$ and need not be covered. Hence the cardioid family, the arc
charts, and the complete exterior family glue on all of $\HH$.
They are explicit, single-valued local constructions with proved open
overlap agreement. This proves the assertion about paths as well as
joint local holomorphy; no inference from numerical interpolation is
made.
\end{proof}

\section{The marked cusp corridor and the sewn boundary value}
In this section $\varepsilon$ is in the sector defined above, and
$\rho=|\varepsilon|$. All unquantified constants are uniform after
decreasing a positive sector radius finitely many times. This radius is
used only for F; the complete-exterior certificate has its independent,
explicit radius in Lemma~\ref{lem:tail}.
Write $u^*(\varepsilon)=f_c^3(0)-1/2$ and
$u_n^*=P_\varepsilon^n(u^*)$. At $\varepsilon=0$,
$u^*(0)=-39/256$, so the selected inverse chains avoid $u=-1/2$.

\begin{lemma}[A marked tilted corridor]\label{lem:marked-cusp}
There exist $r_0,\delta>0$ such that on the smaller sector
\[
 0<|\varepsilon|<r_0,\qquad
 -7\pi/12+1/10<\arg\varepsilon<\pi/12-1/10,
\]
the tilted quotients define one jointly holomorphic family $F_\varepsilon$
on $W_\delta=\{z:\operatorname{dist}(z,[0,1])<\delta\}$, normalised at
$u^*(\varepsilon)$ in the $u$ coordinate. Its mark is the first crossing
of the continued logarithmic orbit, and it is independent of the
admissible line. On the real exterior ray it is the original Kneser
family, expressed in $u$ coordinates.
\end{lemma}
\begin{proof}
We give the orbit and inverse-chart details needed in addition to the
unmarked estimates. At $\varepsilon=0$, put $v_n=-1/u_n^*$.
Then $v_{n+1}=v_n+1+1/(v_n-1)$, so
$v_n=n+O(\log(n+2))$ and
$u_n^*=-1/n+O(\log n/n^2)$. Choose a small $\rho_1>0$ and a fixed
$N_1$ so that $|u_{N_1}^*(0)|<\rho_1/8$. Continuity of this finite
orbit, including its holomorphic moving base, allows the sector radius
to be decreased so that $u_{N_1}^*$ differs from this negative real
value by less than a tenth of its modulus.
With $A=1/\Log(1-\varepsilon)$, start $\chi_{N_1}^*$ at the principal
logarithm and continue by $\chi_{n+1}^*=\chi_n^*+\Delta(\chi_n^*)$.
Initially $Z_{N_1}^*:=A\chi_{N_1}^*=-1/u_{N_1}^*(1+O(\rho/\rho_1))$,
within relative distance $1/5$ of a real number at least $8/\rho_1$.

Besides Lemma~\ref{lem:step}, the exact integral formula gives, for
$|u|<1/8$,
\[
 A\Delta=1+O(|u|+\rho),\qquad
 |\Delta_\chi|\le C\rho^2/|\chi|^2\quad(0<|\chi|\le1).
\]
The last estimate follows by differentiating $u(\chi)$ and using
$|(e^\chi-1)/\chi|\ge1-e^{-1}$ on $|\chi|\le1$.
Here $|u|\le C/|Z|+C\rho$; on $\mathcal G$ it is $O(\rho)$.
Induction therefore increases $\re Z_n^*$ by at least $1/2$ per step
and gives, up to $N_1+10/\rho$ and before any hole,
\[
 Z_n^*=Z_{N_1}^*+(n-N_1)+E_n,\qquad
 |E_n|\le C\log(n+2)+C\rho n.
\]
To close this bootstrap, the resulting $\chi$ orbit is within
$C\rho\log(1/\rho)+C\rho/\rho_1$ of the ray $-\varepsilon\tau$,
$\tau\ge0$, before its first line crossing. The ray's normal coordinate
increases at rate at least $\rho/2$; at its origin it is between
$-\sqrt{\pi^2+1/4}$ and $-1.48$. It crosses by $6.4/\rho$.
Before that crossing it stays at distance at least $1.48$ from the
upper holes and at least $2\pi\sin(5\pi/12)>6$ from the lower holes.
The origin hole is handled by $|\chi|\le1$ and the increasing $\re Z$.
Shrinking $\rho_1$ and then $r_0$ proves that the actual orbit crosses
by $n_0\le N_1+8/\rho$, and that its normal component is strictly
increasing. The injective isotopy $\chi\mapsto\chi+s\Delta(\chi)$
near the crossing puts $\chi_{n_0}^*$ in the tilted strip, including
its chord and excluding its image edge. The three-unit neighbourhoods
in $Z$ of the preceding orbit stay in $|u|<2\rho_1$ and away from
every nonzero hole. This is a first crossing on the logarithmic lift;
other physical visits to a spiralling crescent are not its mark.

For completeness we control the inverse uniformisation at this mark.
The flat model differs from the identity in $Z$ by $O(\rho)$ with
Beltrami norm $O(\rho)$. Put its marked point at $1$ in the exponential
coordinate. For coefficients $t\mu$ with
$|t|<(2\|\mu\|_\infty)^{-1}$ the normalised sphere solutions are
uniformly 3-quasiconformal. Their values on a fixed compact annulus
are uniformly bounded by normality. Holomorphic dependence in $t$
and the Schwarz estimate then bound their displacement at $t=1$ by
$C\|\mu\|_\infty$. The inverse has the same bound on a smaller
annulus, by its boundary degree and the holomorphic inverse theorem.
Taking the logarithm and allowing finitely many adjacent seams gives
\[
 Z(\widetilde F(n_0+z))=Z_{n_0}^*+z+O(\rho)
 \quad (|\re z|\le2,\ |\im z|\le1).
\]
In $Z$, the selected inverse has derivative
$1+O(\rho^2+|Z|^{-2})$. Summing along the noncritical orbit gives
\[
 \sum_{j=N_1}^{n_0}(\rho^2+|Z_j^*|^{-2})\le C\rho+C\rho_1.
\]
The product estimate for these derivatives consequently keeps the
pulled-back window within $1/10$ of $Z_{n_0-k}^*+z$, for
$0\le k\le n_0-N_1$, after decreasing the two radii.
The physical inverse is
\[
 g_\varepsilon(v)=-1/2+\sqrt{v+1/4+\varepsilon^2/4},
\]
with the square root positive on the selected real orbit at
$\varepsilon=0$. All these near-cusp steps lie near $u=0$, so this
branch is jointly holomorphic and noncritical there.

The remaining $N_1$ steps are a fixed finite chain. The parabolic
attracting Abel coordinate has
$\Phi(u)=-1/u+\log(-u)+O(u)$; its inverse near the positive real
orbit-time segment from $u^*(0)$ to $P_0(u^*(0))$ exists because
all these physical points have $w>0$. Fix a small $\delta>0$ in this
parabolic time chart and choose $\rho_1$ small enough that its $O(\rho_1)$
time errors fit within the chart's margin. The compact segment and its
$N_1$ forward images have noncritical inverse neighbourhoods. The window
above, pulled back to $N_1$, differs from that Abel-time window by
$O(\rho_1)+O(\rho)$, since the logarithmic term varies by $O(\rho_1)$
there. Compactness and continuity preserve the finite chain for small
$\rho$. Thus
$F^\ell_\varepsilon(z)=g_\varepsilon^{n_0}
 (\widetilde F(z+n_0))$ is defined on $W_\delta$ with the required
mark and functional equation.

We next verify independence of the line. For two sufficiently close
admissible pairs, at a point of tangential coordinate $t$ their normal
coordinates differ by at most $C\eta_0(1+|t|)+C\rho$, where
$\eta_0$ is the distance between the pairs. Each dynamical step has
normal increase at least $2\rho/5$ and modulus at most $2\rho$.
In at most $C\eta_0(1+|t|)/\rho+C$ forward or backward steps the
first crossing of the other strip is reached. Their total displacement
is at most $C'\eta_0(1+|t|)+C'\rho$; choose $\eta_0$ and $r_0$
so it is smaller than half the distance to any hole. This is possible
because the original distance is at least $c(1+|t|)$.
Fixed iterates give holomorphic chart maps; their jumps are precisely
the new seam. They induce a proper end-preserving quotient map of
winding number one, hence a biholomorphism. Both marks belong to the
same continued base orbit and are its first crossings, so the map
changes their times by exactly $n_0'-n_0$. The inverse chains cancel
that change. Admissible pairs form a connected rectangle, so subdivision
of a path proves equality for arbitrary admissible lines.

Finally freeze an admissible line on a small parameter disc, taking
its angle nonzero and moving $x_0$ so the mark is strictly interior.
Then $n_0$ is locally constant. The real source parametrisation $N$
is holomorphic in $\varepsilon$ at fixed $(t,s)$; relative to one
fixed fibre its Beltrami coefficient is holomorphic in
$L^\infty$ and has norm $<1$ after shrinking the disc. The same
Ahlfors--Bers cancellation as in Theorem~\ref{thm:qc}, followed by
the noncritical inverse chains just proved, gives joint holomorphy.
Line independence patches these parameter discs and the finite time
charts. For $\varepsilon=-i\tau$, $\tau>0$, the horizontal line is
admissible and is precisely the chord between the conjugate fixed
points. Its continued critical orbit first crosses that chord with
the real positive inverse branches. The mark and ordered ends are
the original real Kneser ones. This proves the last assertion.
\end{proof}

\begin{lemma}[The sewn real family and the large cusp collar]
\label{lem:sewn}
For $0<c<1/4$ put $\varepsilon=\sqrt{1-4c}$,
$p=(1-\varepsilon)/2$, $q=(1+\varepsilon)/2$, and
$h=2\pi/|\log(1-\varepsilon)|$.
Let $r_+$ be attracting Koenigs time normalised by $r_+(w_3)=0$,
and let $S$ be the real repelling parametrisation of $(p,q)$,
$S(t+1)=f_c(S(t))$. Choose
\[
 T(t)=r_+(S(t)),\qquad \im T(t)=h/2\quad(t\in\mathbb R).
\]
Gluing $C_+=\{\im r\ge-h/2\}/\mathbb Z$ and
$C_-=\{\im t\le0\}/\mathbb Z$ by $r=T(t)-ih$ gives the marked
sewn sphere and a germ $K_c^W(0)=w_3$. This germ is jointly
holomorphic on a neighbourhood of every real $c\in(0,1/4)$.
For sufficiently small positive $\varepsilon$, $T$ is univalent on
$|\im t|<0.40h-C$; its imaginary part lies above or below $h/2$
according as $\im t$ is positive or negative.
\end{lemma}
\begin{proof}
For every parameter in this real interval, $0<w_3<p<q$.
Every orbit point on $(0,q)$ is noncritical. The attracting Koenigs
function $\sigma$ has $\sigma(w_3)<0$, $\sigma>0$ on $(p,q)$,
and $\sigma'>0$ there, by pulling back its real normalisation near $p$.
The repelling parametrisation is real and strictly decreasing from
$q$ to $p$, with nonzero derivative along $\mathbb R$.
Hence
\[
 T'(t)=\frac{\sigma'(S(t))S'(t)}{\log(1-\varepsilon)\sigma(S(t))}>0,
 \qquad T(t+1)=T(t)+1.
\]
It induces an analytic circle diffeomorphism, with a univalent analytic
collar by compactness of one period. Sewing two punctured discs along
this collar and filling their remote punctures gives a compact genus-zero
surface. Its two ordered ends and the upper time-zero point determine
the uniformisation uniquely. Inverting at that point in the physical
Koenigs chart gives $K_c^W$.

The fixed points, their Koenigs charts, and the noncritical finite orbit
chains depend holomorphically on the parameter near any fixed real $c$.
Thus $T$ does also on a common collar of one period. To see parameter
holomorphy of the normalised sphere without making a boundary-value
assertion, interpolate the small change of this collar map by a fixed
smooth cut-off on its annulus. This produces source maps holomorphic
in the parameter at fixed point, equal to the identity outside the
collar, and with Beltrami coefficients holomorphic in $L^\infty$ and
of norm $<1$. Ahlfors--Bers and the same cancellation calculation as
above give holomorphic coordinates at the mark. The inverse function
theorem proves joint holomorphy of the germ. Different real parameter
discs agree by their uniquely marked sewing construction.

Here are direct quadratic bounds for the larger collar near the cusp.
In the $u$ variable set
\[
 \Psi=A_1\log(u+\varepsilon/2)+A_2\log(u-\varepsilon/2),
 \quad A_1=1/\log(1-\varepsilon),\quad A_2=1/\log(1+\varepsilon).
\]
The residual $E=\Psi\circ P_\varepsilon-\Psi-1$ is
\[
 A_1\Log(1+u-\varepsilon/2)
 +A_2\Log(1+u+\varepsilon/2)-1.
\]
It is holomorphic also at $\varepsilon=0$: the apparent simple poles
cancel. It vanishes at both fixed points. Division by
$u^2-\varepsilon^2/4$, or the holomorphic remainder theorem on a
fixed small disc, gives
$|E(u)|\le C|u+\varepsilon/2||u-\varepsilon/2|$ uniformly.
The exact formulas give $A_1+A_2=1+O(\varepsilon^2)$.

On $0.55<\im\chi<2\pi-0.55$ the forward lifted orbit moves left
by at least $\varepsilon/2$ per step and the backward orbit moves right.
Their total imaginary drift is $O(\varepsilon)$: near either end
$|\im\Delta|\le C\varepsilon^2e^{-|\re\chi|}$, because the limiting
steps are real, and the sum over the middle and the two tails is
$O(\varepsilon)$. Shrinking the radius keeps both orbits at distance
greater than $1/2$ from the holes. Summing $E$ along them gives
\[
 r_+=\Psi+\sum_{k\ge0}E\circ P_\varepsilon^k+\text{constant},\qquad
 r_-=\Psi-\sum_{k\ge1}E\circ P_\varepsilon^{-k}+\text{constant},
\]
where the sums and their $\chi$ derivatives are $O(\varepsilon)$
on every slightly smaller horizontal strip. They are the Koenigs
times, since their limits near the respective fixed points have the
Koenigs logarithmic singularities and satisfy the Abel equation.
Using
$\Psi=A_1\chi+(A_1+A_2)\log(u-\varepsilon/2)$, its second term
has bounded $\chi$ derivative. Thus
$r_\pm'=A_1+O(1)$ uniformly there. Normalising $r_-$ to be real on
$\im\chi=\pi$ gives
\[
 \left|\im r_-+\frac h{2\pi}(\im\chi-\pi)\right|\le C.
\]
The arguments of $1-e^\chi$ are bounded on this horizontal strip,
which proves the uniform constant here for all real parts.

Both $-r_+$ and $-r_-$ have derivative with positive real part on
the convex strip, hence are univalent. Their horizontal boundary
images are proper graphs, tending to opposite real infinities.
Using a smaller strip with boundaries at $0.56$ and $2\pi-0.56$,
the last estimate and the argument principle show that the image of
$r_-$ contains $|\im t|<0.40h-C$.
There $T=r_+\circ r_-^{-1}$, up to its specified period $ih$,
has $T'=1+O(\varepsilon)$. It is univalent as a composition of the
two univalent charts. Since $T(\mathbb R)$ is exactly the line
$\im r=h/2$, integrating $\re T'>0$ along vertical segments proves
the final sign assertion. This is an analytic continuation of the
original circle collar; no hypothesis on a horn coefficient is needed.
\end{proof}

\begin{proposition}[Identification of the upper sewn branch]
\label{prop:sewn-identification}
For sufficiently small real $\varepsilon>0$,
$1/2+F_\varepsilon=K_{1/4-\varepsilon^2/4}^W$ on $W_\delta$.
\end{proposition}
\begin{proof}
Take the admissible line $\chi=\pi i+t e^{i\pi/4}$.
The real continued base orbit has $\chi_n^*<0$ and moves left;
its first mark is in $[-\pi-2\varepsilon,-\pi]$.
Along the strip take a curve from this mark in the positive $t$
direction, keeping its relative transverse position fixed, to the
point $u_m\in(u_1,u_2)$ where $\im\chi=\pi$.
The intervening points $0<\im\chi<\pi$ map to the lower physical
half-plane. Consequently the argument of
$\sigma(u)/\sigma(u^*)$ reaches $+\pi$ at $u_m$. The continued
attracting time $r_+$, with $r_+(\chi_{n_0}^*)=n_0$, has
\[
 \im r_+(u_m)=-h/2.
\]
This equality is exact: the real Koenigs ratio there is negative,
and the path fixes its logarithm. Equivalently the residual estimates
give $-h/2+O(\varepsilon)$, selecting the same member of
$-h/2+h\mathbb Z$. Choose $r_-$ real at $u_m$.
Since the specified transition $T$ has imaginary part $+h/2$ there,
\[
 r_+=T\circ r_--ih
\]
on the common continued chart. This fixes the required branch.

To verify that these charts identify whole quotient surfaces, rather
than only one orbit point, let $U,V$ be the saturated parts of the
tilted strip with respectively
$t<\pi\sqrt2-1$ and $t>-\pi\sqrt2+1$.
The forward orbits on $U$ and backward orbits on $V$ avoid the holes:
in their bounded middle part their imaginary coordinates have the
margins used in Lemma~\ref{lem:sewn}, and the total drift is
$O(\varepsilon)$; at their unbounded ends their real parts keep them
away from all holes. Thus the two Koenigs times and their residual
estimates apply on these charts. Crossing a seam adds exactly one,
since their changes are $O(1)$ and the only possible other logarithmic
periods have modulus $h\to\infty$.
On $U$ and on $U\cap V$, respectively, the argument bounds give
\[
 \left|\im r_++\frac h{2\pi}\im\chi\right|\le C,
 \qquad
 \left|\im r_-+\frac h{2\pi}(\im\chi-\pi)\right|\le C.
\]
On $U\cap V$, $|\im\chi-\pi|\le\pi-1/\sqrt2+O(\varepsilon)$,
so $|\im r_-|<0.39h+C$, within the large collar of
Lemma~\ref{lem:sewn} after shrinking the radius.
On $U\setminus V$ the first estimate gives $\im r_+>-h/2$.
On $V\setminus U$, using
$\arg(1-e^\chi)=\im\chi-\pi+O(1)$ at the positive real end gives
$\im r_-<0$. On the overlap the sign assertion for $T$ says
$\im r_+\ge-h/2$ exactly when $\im r_-\ge0$.

Define a map to the sewn surface by $r_+$ when
$\im r_+\ge-h/2$ and by $r_-$ when $\im r_-\le0$.
These regions cover; on the separating analytic curve their images
are identified by $T-ih$. The map is holomorphic across that curve
in the sewing collar. At the attracting end $\im r_+\to+\infty$,
and at the repelling end $\im r_-\to-\infty$. It is therefore proper
and extends to the two compactified spheres with no additional
preimages of the ends. A seam circuit near the attracting end changes
$r_+$ by exactly one, so its local winding number, and hence its
degree, is one. The map is biholomorphic.
It takes the tilted mark to attracting time $n_0$. Both germs use
attracting time zero at $u^*$ and the same selected finite noncritical
inverse chain from that mark. Their uniquely normalised lifts agree,
proving the assertion on $W_\delta$.
\end{proof}

\begin{theorem}[Quadratic F]\label{thm:F}
The original real exterior Kneser family continues through an upper
punctured neighbourhood of $c=1/4$, crossing the neutral boundary
without an arithmetic restriction on its multipliers. Its boundary
value at every real $0<c<1/4$ is the sewn family $K_c^W$ specified by
the transition $T-ih$. It extends single-valuedly through the whole
upper main cardioid. The normalisation throughout is $F_c(0)=f_c^3(0)$.
\end{theorem}
\begin{proof}
The exact map $c=1/4-\varepsilon^2/4$ is injective on the smaller
sector of Lemma~\ref{lem:marked-cusp}; its image contains an upper
punctured half-disc and adjacent real segments on both sides of the
cusp. That lemma gives joint holomorphy on a common $W_\delta$,
and equality to the real exterior family on the right segment.
Proposition~\ref{prop:sewn-identification} identifies its left segment.
Its upper portion agrees with Theorem~\ref{thm:G}: there is an open
exterior overlap where the horizontal tilted crescent is the straight
one with the same first-orbit marking, and the identity theorem gives
equality on the connected upper cusp neighbourhood.

The multiplier deformation in Lemma~\ref{lem:surgery} is also defined
across the positive real interval $0<\lambda<1$. The map $c(\lambda)$
therefore gives parameter germs across every $0<c<1/4$ for the same
global upper family. The sewn family has local holomorphic germs
along that entire interval by Lemma~\ref{lem:sewn}. They agree near
the cusp; a chain of overlapping discs along any compact subinterval
propagates equality by the identity theorem. Their upper boundary
values coincide with $K_c^W$ at every real point of the interval.
The full upper-cardioid assertion follows from the already glued
family of Theorem~\ref{thm:G}. All claims concern regular time germs
and their specified continuation charts.
\end{proof}

\section{Scope of the continuation and certificate}
The cubic $c^3+2c^2+c+1$ has one root
\[
 c_{\rm per}=-0.1225611668766536\ldots+0.7448617666197442\ldots i
\]
in $\HH$. Its critical point has exact period $3$, because
$f_c^3(0)=c(c^3+2c^2+c+1)$ and
$c_{\rm per},c_{\rm per}^2+c_{\rm per}\ne0$.
This prevents a single global dynamical quasiconformal conjugacy from an
escaping exterior real polynomial to every upper parameter: a conjugacy
must preserve the unique critical point and whether its orbit is periodic.
It does not refute G as a statement about height germs.
In fact $c_{\rm per}\in\mathcal O$, so Theorem~\ref{thm:outer} handles this
parameter using the critical-value mark and a forward push, rather than
pulling a locally univalent mark back through a critical point.
At $c_{\rm per}$ its time-zero derivative is
$4c_{\rm per}(c_{\rm per}^2+c_{\rm per})U'_{c_{\rm per}}(0)\ne0$,
while its forward image at time $1$ has zero derivative. Germ holomorphy must not be confused with
univalence at all translated times.

The proof of G uses Theorem~\ref{thm:straight}, the analytic and certified
entry lemmas, Lemma~\ref{lem:entry-patch}, and the explicit transport and
open-overlap comparison. The proof of F additionally uses the marked
tilted corridor and the exact $T-ih$ sewing identification. In particular
the mark is never replaced by an arbitrary interior section or a Fourier
constant mode without comparison.

The finite certificates have separate roles. The geometry certificate
checks the explicit rational constants in the cusp estimates; the entry
certificate proves a full compact parameter cover and the preceding
noncritical orbit predicates. Its independent verifier proves exact
coverage as well as replaying every interval inequality. Neither is a
sampling test. The analytic tail and the $R>1/2$ formula cover the
noncompact remaining ranges. To reproduce the finite checks from the
repository root run
\begin{verbatim}
python3 docs/certify_quad_geometry.py \
  --output docs/proofs/quadratic-geometry-certificate.json
python3 docs/certify_quad_entry.py \
  --output docs/proofs/quadratic-entry-core.json
python3 docs/verify_quad_entry.py \
  docs/proofs/quadratic-entry-core.json \
  --output docs/proofs/quadratic-entry-verification.json
\end{verbatim}
The proofs do not assert continuation at the cusp itself, at the real
endpoint $-3/4$, or on a common slit height plane for all parameters.

\section{The present theta/Taylor iteration cannot converge on all of $\HH$}
\label{sec:numerical-obstruction}
Here ``the present iteration'' means \texttt{docs/quad\_kneser.py}:
the two fixed-point inverse Schr\"oder series, their Abel inverse branches,
one-sided Fourier reconstruction, the functional-equation band, the
unit-circle Cauchy projection, and replacement of its constant Taylor
coefficient by $w_3$. Its whole-parameter convergence is a different
claim from Theorems~\ref{thm:F} and~\ref{thm:G}.

\begin{theorem}[Exact obstruction to the requested global convergence]
\label{thm:theta-no-global}
The present Schr\"oder-based theta/Taylor scheme cannot be defined, and
hence cannot converge to the specified family, at every $c\in\HH$.
Removing its neutral-band refusal, increasing precision or Taylor order,
moving its time Taylor centre, or fixing the constant mode of theta does
not repair this obstruction while the same two Schr\"oder charts are used.
An explicit counterexample is $c_{\rm res}=1/4+i/2$.
\end{theorem}
\begin{proof}
Its upper fixed point is $L=i/2$ and its multiplier is $\lambda=i$.
In $v=w-L$ the exact map is $v\mapsto iv+v^2$.
The upper \texttt{Side} object requires
\[
 u(s)=s+\sum_{k\ge2}u_ks^k,\qquad
 iu(s)+u(s)^2=u(is).
\]
Coefficient comparison gives precisely the recurrence used in
\texttt{schroeder\_inverse}:
\[
 (i^k-i)u_k=\sum_{j=1}^{k-1}u_ju_{k-j},\qquad u_1=1.
\]
At the first three nonsingular orders it forces
\[
 u_2=\frac{-1+i}{2},\qquad
 u_3=\frac{-1-i}{2},\qquad
 u_4=\frac{1-5i}{4}.
\]
At order five it then forces
\[
 0=(i^5-i)u_5=2u_4+2u_2u_3=\frac{3-5i}{2}\ne0.
\]
Thus even a formal inverse Schr\"oder series with nonzero linear
coefficient does not exist. For a general linear coefficient $A\ne0$
the forced preceding coefficients are $A^ku_k$ and the contradiction
becomes $0=A^5(3-5i)/2$; changing the scaling does not help.
A time translation changes the Schr\"oder argument by a nonzero
factor $e^{z_0\Log\lambda}$, and the constant-mode gauge is also an
Abel-time translation. Neither can turn a nonexistent invertible
Schr\"oder germ into an existing one. This obstruction is internal to
the required upper chart, independent of the downstream Taylor circle.

As written, \texttt{build} already rejects this point via
$\bigl||\lambda|-1\bigr|<0.02$. If that guard is simply removed,
the series construction encounters the displayed nonzero numerator
and zero denominator. In addition, its depth formula divides by
$|\log|\lambda||=0$. The series contradiction proves that this is
not just an avoidable guard or depth-estimation failure.
The accompanying \texttt{certify\_quad\_theta\_obstruction.py} checks
all of the displayed coefficients and the parameter using exact
fractions in $\mathbb Q(i)$, with output
\texttt{quadratic-theta-obstruction.json}.
\end{proof}

\begin{proposition}[Dense neutral obstructions and loss of a uniform chart]
\label{prop:theta-resonances}
The obstruction to an invertible Schr\"oder chart occurs at every
root-of-unity upper multiplier, a dense set on the open upper neutral
arc. Already near $\lambda=i$ its fifth coefficient has a genuine pole.
\end{proposition}
\begin{proof}
Let $\lambda=e^{2\pi i p/q}$ with $0<p/q<1/2$ and $(p,q)=1$.
Then $c=\lambda/2-\lambda^2/4$ has
$\im c=\sin(2\pi p/q)(1-\cos(2\pi p/q))/2>0$.
If an invertible formal series $u$ conjugated
$g(v)=\lambda v+v^2$ to $s\mapsto\lambda s$, then
$g^{\circ q}\circ u=u$, hence $g^{\circ q}$ would be the identity
as a formal series. It is a polynomial of degree $2^q$, so this is
impossible. Density follows from density of the rational angles.

For $\lambda$ near $i$, $u_2,u_3,u_4$ are analytic rational functions
of $\lambda$, since their denominators are nonzero at $i$. Therefore
\[
 u_5(\lambda)=\frac{2u_4(\lambda)+2u_2(\lambda)u_3(\lambda)}
 {\lambda^5-\lambda},\qquad
 \lim_{\lambda\to i}(\lambda-i)u_5(\lambda)=\frac{3-5i}{8}.
\]
In particular the normalised coefficients are unbounded in any
punctured neighbourhood of this resonance. A uniform analytic radius
\emph{and uniform modulus bound} for these charts there would give a
uniform fifth-coefficient bound by Cauchy's estimate, a contradiction.
This conclusion does not assert a pole of the marked family $F_c$:
it is a pole of this particular linearisation representation.
It does not exclude obtaining a neutral value by a separately proved
parameter limit of assembled outputs. Such a limit is an additional
procedure, not evaluation or convergence of the present per-parameter
Schr\"oder series at that neutral point.
The general resonance phenomenon is discussed in~\cite{BMS}; the
quadratic counterexample here is an exact coefficient calculation.
\end{proof}

\begin{proposition}[Fixed-degree Taylor iteration is not an exact limit]
\label{prop:finite-taylor}
At a fixed finite \texttt{nt}, convergence of the stored Taylor
coefficient vector is not convergence to an exact nonconstant solution
of $F(z+1)=F(z)^2+c$ on an open height overlap. For example, no such
limit with the prescribed normalisation can exist at $c=i$.
\end{proposition}
\begin{proof}
A convergent finite coefficient vector defines a polynomial $P$.
If $P(z+1)=P(z)^2+c$ on any nonempty open set, the identity theorem
gives a polynomial identity. For a nonconstant polynomial its two
sides have degrees $d$ and $2d$, impossible. A constant solution
must be a fixed point. At $c=i$, the critical orbit gives $w_3=-i$,
but $f_i(-i)=-1+i\ne-i$, so even the constant possibility is excluded.
These last two equalities are also checked by the exact certificate.
The observation concerns an exact limit at fixed truncation. It does
not rule out accurate finite approximations or a proved limit with
all discretisation and truncation errors tending to zero.
\end{proof}

\subsection*{A valid local theorem, with its hypotheses exposed}
There is a useful convergence statement once regular charts, a function
space, contraction, consistency, and branch identity have been verified.
It is stated here to identify the actual required inequalities; no
quadratic parameter-wide contraction certificate is claimed.

Fix a nonneutral parameter for which the following boundary construction
at the marked time zero is legitimate. Take the unit circle centred at zero,
sampling heights $0<\delta_s<\delta_a<1$, and choose
\[
 \max\left(\sqrt{1/4+\delta_s^2},
  2\sin\bigl(\tfrac12\arcsin\delta_a\bigr)\right)<r<1.
\]
On the complex Wiener space
$A_r=\{\sum a_kz^k:\sum|a_k|r^k<\infty\}$,
restrict to the affine normalisation $f(0)=w_3$.
Assume the two selected regular maps $S_\pm$, their Abel inverse
branches $A_\pm$, and the selected quadratic inverse in the left band
are analytic on neighbourhoods of \emph{all} values attained by a
closed $A_r$ ball $B(p,R)$. Branch and period choices are fixed on
this ball. Take the exact continuous Fourier integrals of
$A_\pm(f(t\pm i\delta_s))-(t\pm i\delta_s)$, retain the appropriate
one-sided modes, and reconstruct on the arcs
$\pm\im z\ge\delta_a$. The two band values are
$f_c(f(z-1))$ and the chosen inverse of $f(z+1)$.
Denote this boundary data by $B_c(f)$, its Cauchy projection by
$\mathcal P$, and define
\[
 \mathcal T_c(f)=w_3+\mathcal P B_c(f)
                         -(\mathcal P B_c(f))(0).
\]
Separated heights make the one-sided mode sums absolutely convergent
even for nonperiodic candidate inputs. This is a modified ideal operator,
not the finite same-height program silently relabelled.

\begin{theorem}[Conditional local convergence and discretisation error]
\label{thm:theta-local}
Assume this full-ball analytic construction has been verified and
\[
 \sup_{f\in B(p,R)}\|D\mathcal T_c(f)\|\le q<1,
 \qquad d:=\|\mathcal T_c(p)-p\|\le(1-q)R.
\]
Then its exact iterations starting anywhere in $B(p,R)$ converge to
its unique fixed point $f_*$. Suppose computed iterates remain in the
ball and obey
\[
 f_{n+1}=\mathcal T_c(f_n)+e_n,\qquad \|e_n\|\le\eta_n.
\]
Then
\[
 \|f_n-f_*\|\le q^n\|f_0-f_*\|
       +\sum_{j=0}^{n-1}q^{n-1-j}\eta_j.
\]
If $\eta_n\to0$, they converge to $f_*$. If the marked family has
independently been proved to belong to this ball and to satisfy
$\mathcal T_c(F_c)=F_c$, then $f_*=F_c$.
\end{theorem}
\begin{proof}
The derivative bound integrated along the line segment between two
ball elements proves the $q$-Lipschitz estimate. For $f$ in the ball,
$\|\mathcal T_c(f)-p\|\le qR+d\le R$. The closed ball is complete,
so Banach's theorem proves the first assertion. The inexact recursion
gives $\|f_{n+1}-f_*\|\le q\|f_n-f_*\|+\eta_n$; induction gives
the displayed error formula. For $\eta_n\to0$, split that sum at a
fixed index: its finite initial part tends to zero geometrically and
the remaining part is at most $\sup_{j\ge J}\eta_j/(1-q)$.
For a uniform margin $\eta_n\le(1-q)R-d$, ball membership itself
also follows by induction. Finally the independently identified
$F_c$ is another fixed point in this ball, so uniqueness identifies it.
\end{proof}

Each numerical error $\eta_n$ here must include the regular-map
approximation, its inverse and branch error, Fourier truncation and
quadrature, circle quadrature, Taylor truncation, and roundoff.
For example, writing $\gamma=e^{-2\pi(\delta_a-\delta_s)}<1$,
a sampling-line bound $B_\theta$ gives the mode tail
$B_\theta\gamma^M/(1-\gamma)$; a boundary-data bound $B_0$ gives
the output $A_r$ Taylor tail $B_0r^N/(1-r)$.
These follow respectively from bounding each Fourier coefficient by
$B_\theta$ and each Cauchy coefficient by $B_0$.
They do not control quadrature or branch errors by themselves.
At fixed precision and fixed truncations the theorem supplies an error
floor bound, rather than exact convergence to the marked family.

Neither Theorem~\ref{thm:G} nor the 4,013 entry boxes verify these
derivative or chart hypotheses: their predicates concern parameter
holomorphy and critical-orbit geometry, not the numerical iteration.
Moreover Theorem~\ref{thm:theta-no-global} makes the regular-chart
hypothesis impossible on all of $\HH$ for the present two-sided
Schr\"oder scheme. A genuinely whole-half-plane numerical method must
replace its neutral upper-end representation, for example by quotient
uniformisation charts; centre or constant-mode changes alone do not
suffice. Such a replacement would require its own convergence proof.

\section{A general theory of marked cylinder continuation}
\label{sec:general-cylinders}
The transferable object is a cylinder equipped with ordered ends, a
physical orbit marking, and specified analytic seam charts. It is not
the formula $w^2+c$. This section proves abstract continuation and
invariance statements. Membership in the class below is a geometric
hypothesis to be verified for each family; it is not a consequence of
having a simple saddle-node or a single critical point.

\subsection{An admissible class and a continuation criterion}
Let $D$ be a connected parameter domain, $f_\nu$ a jointly holomorphic
family on the physical regions in use, and $s(\nu)$ a holomorphic marked
section. Fix $N\ge1$ and put $m_N(\nu)=f_\nu^N(s(\nu))$. Call a family
\emph{admissible for marked cylinder continuation} on $D$ if it has an
open cover $D_i$ with the following data.
\begin{enumerate}
\item Each $D_i$ carries a holomorphic family of seam quotients $X_i$,
  with ordered ends, biholomorphic fibrewise to $\C^*$. The physical seam
  map is $f_\nu$. There are jointly holomorphic uniformisations $Q_i$;
  only multiplication by a nonzero parameter function is allowed when
  changing their normalisation. A locally uniform Beltrami construction
  as in Theorem~\ref{thm:qc} is one sufficient way to obtain these data.
\item A finite orbit point $m_i=f_\nu^{n_i}(s)$ is in a physical interior
  or seam chart of $X_i$. A specified finite chain of local biholomorphic
  forward or inverse dynamical charts transports its neighbourhood to a
  neighbourhood of $m_N$. Denote its composition by $R_i$; thus
  $R_i(m_i)=m_N$. All inverse branches in this chain are specified and
  noncritical, and the chain varies holomorphically in $\nu$.
\item On overlaps the quotient identifications are jointly holomorphic
  isomorphisms preserving the ordered ends and the orbit marking. Their
  physical chart identifications are compositions of analytic seam
  steps. These steps have the recorded integer time displacement and
  agree with the chosen transport chains to $m_N$. In particular, a
  later return to a quotient is not identified with a first entry merely
  because the two points have the same quotient coordinate.
\end{enumerate}
The third condition is a diagram of physical charts, rather than an
assumption that the desired height germs already agree. It can be checked
by comparing analytic seam maps and the finite inverse branches, as in
Lemma~\ref{lem:entry-patch}. If two cylinders are obtained by different
constructions, an isomorphism on an open parameter overlap must first
be proved, as in Lemma~\ref{lem:transport-compare}.

\begin{theorem}[Marked cylinder continuation criterion]
\label{thm:atlas-criterion}
An admissible family determines a single-valued jointly holomorphic
family of locally univalent height germs $F_\nu$ at zero, with
$F_\nu(0)=m_N(\nu)$. It satisfies
\[
 F_\nu(z+1)=f_\nu(F_\nu(z))
\]
in its continued dynamical charts. It is independent of the listed
quotient coordinate normalisations and of admissible changes of entry
index. A common global height domain is not implied.
\end{theorem}
\begin{proof}
In a physical chart about $m_i$, set
\[
 A_i(w)=\frac{1}{2\pi i}\log\frac{Q_i([w])}{Q_i([m_i])},\qquad A_i(m_i)=0.
\]
Here $[w]$ denotes the seam quotient, and the logarithm is the local
branch with value zero at the mark. Its derivative is nonzero. The
orientation is chosen so that continuation across one forward seam
changes $A_i$ by $+1$. Let $U_i$ be its local inverse at zero and define
$F_i=R_i\circ U_i$. Both are jointly holomorphic by the holomorphic
inverse theorem; the finite chain gives the functional equation in
successive charts. Notice that $R_i$ is a dynamical transport, so no
unproved evaluation of $U_i$ at the large height $N-n_i$ is required.

An automorphism of $\C^*$ preserving its two ordered ends is $q\mapsto
\alpha q$. The scale $\alpha$ cancels in the displayed ratio. On an
overlap, the quotient isomorphism therefore identifies the local
logarithms at the corresponding marked classes. The specified physical
seam chain identifies their inverse germs after transport to $m_N$.
The recorded integer seam displacement removes the possible ambiguity
in the lift. Consequently $F_i=F_j$ as physical germs. This also proves
consistency on triple overlaps and independence of an admissible entry
index. Gluing gives a family on $D$, even if $D$ is not simply connected.
Noncriticality of the transport chains proves local univalence. Nothing
in this construction supplies a height disc uniform over noncompact $D$.
\end{proof}

\begin{proposition}[Quantitative stability of finite entry witnesses]
\label{prop:finite-witness-stability}
Let $K\subset D$ be compact and consider finitely many orbit witnesses
with indices at most $M$. Suppose their base orbit points have closed
physical discs of radius $r>0$ on which $f_\nu$ is holomorphic,
$|f_\nu'|\le L$, and $|f_\nu'|\ge a>0$ on every disc used for an inverse
step. Suppose every required real entry predicate $\Psi$ has margin
$\Psi\ge\sigma>0$ at its witnessed point and is $B$-Lipschitz on the
corresponding discs. The constants are uniform on $K$ and over the
finite witnesses.

Let $\widetilde f$ have the same initial section, be holomorphic on
these discs and satisfy $|\widetilde f-f|\le e$ there. Let the corresponding
perturbed predicates differ by at most $\eta$ on the discs. Put
$S_M=\sum_{j=0}^{M-1}L^j$. If
\[
 eS_M<r/2,\qquad BeS_M+\eta<\sigma,\qquad 2e/r<a,
\]
then all these finite entry inequalities and inverse-step
noncriticality predicates persist for the perturbed witnesses.
This assertion concerns finite orbit data; it does not construct the
perturbed global quotient or its infinite ends.
\end{proposition}
\begin{proof}
For the $j$th orbit pair let $E_j$ be its physical distance. Starting
with $E_0=0$, the derivative bound along the segment in the base disc
gives $E_{j+1}\le LE_j+e$. Induction gives
$E_j\le e\sum_{k=0}^{j-1}L^k<r/2$, so every estimate is evaluated in
its stated disc. At a witnessed point the perturbed predicate is at
least $\sigma-BE_j-\eta>0$. For an inverse step the perturbed orbit
point has a disc of radius $r/2$ within the base disc. Cauchy's estimate
applied to $\widetilde f-f$ gives derivative difference at most $2e/r$.
The derivative lower bound therefore leaves
$|\widetilde f'|\ge a-2e/r>0$.
\end{proof}

\begin{proposition}[Propagation of a sewn identification]
\label{prop:general-sewn}
Suppose an admissible family extends to parameter neighbourhoods of a
connected real interval $I$, and a prescribed sewn family has jointly
holomorphic regular germs on the same neighbourhoods. If they agree on
one nonempty open parameter overlap, they agree along all of $I$ that
is connected to this overlap by these neighbourhoods.
\end{proposition}
\begin{proof}
Shrink the parameter discs so that their height germs are represented
on a common local height disc on each adjacent overlap. The identity
theorem gives equality there and then in the next parameter disc.
A finite chain covers the path to any fixed point of $I$. Applying the
argument to all such paths proves the assertion.
\end{proof}
This is a general form of the propagation step in F. The initial
cusp-to-sewing identification is an additional theorem, not supplied
by abstract gluing. Likewise a compact interval certificate supplies
only the entry and regularity predicates it actually verifies.

\subsection{A concrete class beyond degree two}
\begin{proposition}[Unicritical polynomial saddle-nodes]
\label{prop:unicritical-fold}
For every integer $d\ge2$, let
\[
 f_c(w)=w^d+c,\quad x_d=d^{-1/(d-1)},\quad
 c_d=\frac{d-1}{d}x_d,\quad \epsilon=c_d-c.
\]
At $(c_d,x_d)$ the family has a simple real saddle-node. In $u=w-x_d$
its expansion is
\[
\begin{aligned}
 f_{c_d-\epsilon}(x_d+u)-x_d
 &=u-\epsilon+a_2u^2+a_3u^3+O(u^4),\\
 a_2&=\frac{d(d-1)}2x_d^{d-2},\qquad
 a_3=\frac{d(d-1)(d-2)}6x_d^{d-3}.
\end{aligned}
\]
Its iterative residue is
\[
 \rho_d=1-\frac{a_3}{a_2^2}=\frac{d+1}{3(d-1)}.
\]
For small real $\epsilon>0$, the positive fixed points $L_-<x_d<L_+$
have multipliers
\[
 \lambda_\pm=1\pm2\sqrt{a_2\epsilon}+O(\epsilon).
\]
The critical orbit from zero increases to $L_-$, and every orbit point
after zero is positive and noncritical. For $\epsilon<0$ it increases
to infinity. At $\epsilon=0$ it increases to $x_d$.
\end{proposition}
\begin{proof}
The identities $dx_d^{d-1}=1$ and $x_d^d=x_d/d$ prove the fixed-point
and multiplier assertions at the cusp. Taylor expansion gives $a_2,a_3$,
and $x_d^{1-d}=d$ gives the residue formula. The fixed-point equation
is $-\epsilon+a_2u^2+O(u^3)=0$, hence
$u_\pm=\pm\sqrt{\epsilon/a_2}+O(\epsilon)$ and the claimed multipliers.
For $c>0$ the map is increasing on the nonnegative real axis.
The function $x^d+c-x$ has its unique minimum there at $x_d$,
with value $-\epsilon$. If $\epsilon>0$ is small, it is positive on
$[0,L_-)$ and the map sends this interval into $[0,L_-)$; its orbit
therefore increases to $L_-$. At zero minimum the same argument gives
limit $x_d$. At positive minimum the orbit increases and cannot have
a finite limit, since there is no nonnegative fixed point.
Finally $f_c'(w)=dw^{d-1}>0$ at every positive orbit point.
\end{proof}
This proposition proves the local fold and real marking hypotheses
for an infinite class. It does not prove its global cylinder
admissibility. For $d>2$ there are additional fixed points, and the
quadratic collision identity $\zeta+\eta=-1/a$ has no direct substitute.
Global injectivity, the chosen two ends, neutral crossings, orbit entry,
and transport comparisons must be checked afresh. Near-parabolic
renormalisation for unicritical maps of all degrees is already developed
in \cite{Cheritat}; the claim here is the explicit marked-continuation
criterion, not priority for that existing theory.

\subsection{Coordinate invariants and their exact scope}
Let a simple parabolic germ have attracting and repelling Fatou
coordinates $\phi_{\rm att}$ and $\phi_{\rm rep}$. Its two horn lifts,
on their respective ends, have convergent expansions
\[
 H^+(t)=t+\beta_++\sum_{n\ge1}h_n^+e^{2\pi int},\qquad
 H^-(t)=t+\beta_-+\sum_{n\ge1}h_n^-e^{-2\pi int}.
\]
The coordinates here refer to the same repelling and attracting charts
at both ends. Independent changes at each end are not allowed.

\begin{proposition}[Phase-free horn coordinates]
\label{prop:horn-invariants}
Under the coordinate translations
$\widetilde\phi_{\rm att}=\phi_{\rm att}+a$ and
$\widetilde\phi_{\rm rep}=\phi_{\rm rep}+b$, one has
\[
 \widetilde\beta_\pm=\beta_\pm+a-b,\quad
 \widetilde h_n^+=h_n^+e^{-2\pi inb},\quad
 \widetilde h_n^-=h_n^-e^{2\pi inb}.
\]
If $h_1^+\ne0$, the quantities
\[
 \delta=\beta_--\beta_+,\qquad
 J_n^+=\frac{h_n^+}{(h_1^+)^n}\ (n\ge2),\qquad
 J_n^-=h_n^-(h_1^+)^n\ (n\ge1)
\]
are invariant under arbitrary complex coordinate translations. They
determine the horn pair up to these translations. Together with the
formal parabolic data, equality of these quantities gives local analytic
conjugacy by the existing \'Ecalle--Voronin classification theorem.
If $h_1^+=0$, its vanishing is invariant, and the cross-products
$(h_n^+)^m(h_m^-)^n$ and same-end ratios
$(h_n^+)^k/(h_k^+)^n$, where defined, remain invariant.
\end{proposition}
\begin{proof}
The new horn lifts are $\widetilde H^\pm(t)=H^\pm(t-b)+a$.
Substitution gives all three transformation rules and immediately the
invariance assertions. Choose $b$ with $e^{2\pi ib}=h_1^+$ and then
$a=b-\beta_+$. The normalised upper lift has constant zero and first
coefficient one; its remaining coefficients are $J_n^+$.
The lower lift has constant $\delta$ and coefficients $J_n^-$.
The choice of logarithm changes $b$ by an integer, which changes none
of these normalised lifts. Equality of the listed coefficients gives
equality of convergent horn germs, not merely of formal expansions.
The local classification assertion is precisely the classical theorem
recalled in \cite{DS}, applied to these normalised horn pairs.
\end{proof}
These are explicit coordinates on existing horn-map moduli, not newly
discovered replacements for \'Ecalle--Voronin invariants. Individual
$|h_n^\pm|$ are invariant under real repelling-time shifts, but not under
arbitrary complex shifts. The ratios above avoid that restriction.
With the standard matched logarithm branches used in \cite{DS},
$\delta=-2\pi i\rho$; this constant difference is then part of the
formal data, not an additional independent modulus.

For a one-parameter real saddle-node unfolding, the multiplier expression
\[
 p=-\Log\lambda_-\,\Log\lambda_+
\]
is invariant under analytic spatial conjugacy and parameter change,
using the logarithms tending to zero and the fixed ordering. In the
unicritical example $p=4a_2\epsilon+O(\epsilon^2)$, so it is a local
parameter. This product is analytic in $\epsilon$: on the square-root
parameter the two multiplier logarithms are exchanged by the sign
change, leaving their product unchanged.

\begin{proposition}[An intrinsic spectral deformation invariant]
\label{prop:spectral-deformation}
Let a simple saddle-node unfolding, in coordinates at its cusp, be
\[
 g_\epsilon(u)=u-\gamma\epsilon+a_2u^2+a_3u^3
       +b_1\epsilon u+O(u^4,\epsilon u^2,\epsilon^2),
 \qquad a_2\gamma\ne0.
\]
For its two nearby fixed-point multipliers choose logarithms tending
to zero, and define
\[
 p=-\ell_-\ell_+,\quad \ell_\pm=\Log\lambda_\pm,\qquad
 \mathcal R(p)=-\frac{\ell_-+\ell_+}{p}
             =\frac1{\ell_-}+\frac1{\ell_+}.
\]
Then $p$ is an analytic local parameter and $\mathcal R$ extends
holomorphically to $p=0$, with $\mathcal R(0)=1-a_3/a_2^2$.
The function $\mathcal R(p)$, including every Taylor coefficient, is
invariant under analytic conjugacy of unfoldings and analytic parameter
change with nonzero derivative. This is a spectral invariant; it is
not claimed to classify unfoldings completely.
\end{proposition}
\begin{proof}
Set $\epsilon=t^2$ and $A=\sqrt{\gamma/a_2}$. The fixed points are
\[
 u_\pm=\pm At-\left(\frac{a_3\gamma}{2a_2^2}
                     +\frac{b_1}{2a_2}\right)t^2+O(t^3).
\]
Differentiation gives
$\lambda_\pm=1\pm2\sqrt{a_2\gamma}\,t+
2a_3\gamma t^2/a_2+O(t^3)$; the $b_1$ terms cancel.
The involution $t\mapsto-t$ exchanges the two points. Hence the
following symmetric expressions are analytic in $\epsilon$:
\[
 p=4a_2\gamma\epsilon+O(\epsilon^2),\qquad
 \ell_-+\ell_+=-4a_2\gamma
                  (1-a_3/a_2^2)\epsilon+O(\epsilon^2).
\]
Division and the inverse function theorem prove the first claims.
Analytic spatial conjugacy preserves each multiplier; reparametrisation
composes both symmetric expressions with the same parameter change.
Expressing their ratio in the intrinsic coordinate $p$ removes that
change and proves invariance.
\end{proof}

\begin{corollary}[Classification of the fixed-point spectral data]
\label{cor:spectral-classification}
Within these simple unfoldings, $\mathcal R(p)$ is a complete invariant
of the unordered pair of nearby fixed-point multiplier functions, up
to analytic parameter change. It is not a complete invariant of the
maps themselves.
\end{corollary}
\begin{proof}
The two logarithms are exactly the roots of
\[
 X^2+p\mathcal R(p)X-p=0.
\]
Thus $\mathcal R$ reconstructs the unordered logarithms and hence the
multiplier pair. Conversely their sum and product reconstruct $p$ and
$\mathcal R$. If two such functions agree, composing the inverse
intrinsic parameter of one unfolding with that of the other supplies
the required analytic reparametrisation of their spectral data. No
conjugacy of their physical maps follows from this reconstruction.
\end{proof}

\begin{proposition}[Identical cusp germ, inequivalent unfoldings]
\label{prop:same-cusp-different-unfolding}
For $\kappa\in\C$, consider the polynomial unfoldings
\[
 g_{\epsilon,\kappa}(u)=u-\epsilon+u^2+\kappa\epsilon u^3.
\]
All have the identical cusp germ $g_{0,\kappa}(u)=u+u^2$, including
the same parabolic horn-map moduli. Nevertheless
\[
 \mathcal R_\kappa(p)
   =1+\left(\frac1{12}-\frac\kappa4\right)p+O(p^2).
\]
Consequently distinct $\kappa$ give inequivalent analytic unfoldings
under spatial conjugacy and analytic reparametrisation at the cusp.
\end{proposition}
\begin{proof}
In the square-root parameter $\epsilon=t^2$, the two roots are
$u_\pm=\pm t-\kappa t^4/2+O(t^7)$. Thus
\[
 \lambda_\pm=1\pm2t+2\kappa t^4+O(t^7),\quad
 \ell_-+\ell_+=-4t^2+(4\kappa-8)t^4+O(t^6),\quad
 p=4t^2+\frac{20}{3}t^4+O(t^6).
\]
Taking the ratio gives
$\mathcal R=1+(1/3-\kappa)t^2+O(t^4)$, which is the asserted
formula in $p$. Its linear coefficient is injective in $\kappa$;
Proposition~\ref{prop:spectral-deformation} forbids an equivalence
between two distinct values. Equality of the cusp germs proves the
horn-moduli assertion directly, without a numerical comparison.
\end{proof}
This gives a verified distinction between pointwise parabolic data
and deformation data. Its novelty relative to existing saddle-node
unfolding moduli is not asserted: multiplier functions are classical
conjugacy invariants, and $\mathcal R$ is their intrinsic combination.

Whenever cylinder transition maps in an unfolding have the horn form
above and possess sectorial asymptotic expansions
$J_n^\pm(p)\sim\sum_{r\ge0}j_{n,r}^\pm p^r$, their coefficients are
invariants of that unfolding. The assertion follows from the exact
translation rules and uniqueness of asymptotic coefficients.
Existence of such transition maps and expansions is an additional
hypothesis, and it does not follow from the pointwise parabolic
classification. Whether these deformation coefficients provide
information beyond existing unfolding moduli is a literature and
classification question, not established here.

\begin{proposition}[Relative orbit marks and their collision locus]
\label{prop:relative-marks}
Suppose an admissible quotient also carries $k$ labelled holomorphic
sections $v_j$, arising for example from separately tracked singular
orbits. With distinguished mark $m$, set
\[
 \Xi_j(\nu)=\frac{Q_\nu(v_j(\nu))}{Q_\nu(m(\nu))}\in\C^*.
\]
These are holomorphic and invariant under isomorphisms preserving the
ordered ends and all labels. So is the collision equation
\[
 \Delta(\nu)=\prod_j(\Xi_j-1)\prod_{i<j}(\Xi_i-\Xi_j)=0.
\]
On $D\setminus\{\Delta=0\}$, continuation of the labelled tuple defines
a homomorphism into the fundamental group of the ordered configuration
space of $k$ points in $\C\setminus\{0,1\}$, well-defined up to base-point
conjugacy. This braid data is preserved by labelled conjugacy of such
families. It is not claimed to be a complete global classification.
\end{proposition}
\begin{proof}
The multiplier in a coordinate change $Q\mapsto\alpha Q$ cancels in
every ratio, proving both holomorphy and invariance. Noncollision makes
$\nu\mapsto(\Xi_1,\ldots,\Xi_k)$ a map into the indicated configuration
space. Its induced map on fundamental groups proves the last assertion.
Changing the initial configuration identifies the fundamental groups
by a path, giving exactly the stated conjugacy ambiguity.
\end{proof}
Two successive points of the same orbit represent the same quotient
class and hence give $\Xi=1$; counting them as independent marks would
produce no invariant. The relative-mark problem therefore becomes
substantive for distinct tracked orbits. A unicritical family by itself
need not furnish these extra marks. Also, a collision in this quotient
does not by itself prove a singularity of the physical height germ.

\subsection{A general obstruction at roots of unity}
\begin{theorem}[Finite-order criterion and polynomial obstruction]
\label{thm:general-resonance}
Let $g(u)=\lambda u+O(u^2)$, with $\lambda$ a primitive $q$th root of
unity. It is formally conjugate to $u\mapsto\lambda u$ if and only if
$g^{\circ q}(u)=u$ as a formal series. If it is a holomorphic germ, this
condition also gives a holomorphic conjugacy. If $g^{\circ q}\ne\mathrm{id}$
and its first nonlinear term is $Au^r$, then $r\equiv1\pmod q$.
In particular a polynomial of degree $d\ge2$ at any fixed point with
root-of-unity multiplier has no invertible formal Schr\"oder coordinate.
\end{theorem}
\begin{proof}
Conjugacy to the linear rotation makes the $q$th iterate the identity.
Conversely, if $g^{\circ q}=\mathrm{id}$, the average
\[
 L(u)=\frac1q\sum_{j=0}^{q-1}\lambda^{-j}g^{\circ j}(u)
\]
has derivative one and obeys $L\circ g=\lambda L$ by shifting the
summation index. It is formally invertible, and is holomorphic when
$g$ is. For the next assertion put $h=g^{\circ q}=u+Au^r+O(u^{r+1})$.
The identity $h\circ g=g\circ h$ at order $r$ gives
$A\lambda^r=\lambda A$. Hence $\lambda^{r-1}=1$.
For a polynomial $P$ of degree $d\ge2$, formal identity of its $q$th
iterate near a fixed point would be a polynomial identity. But
$\deg(P^{\circ q})=d^q>1$, a contradiction.
\end{proof}
Thus replacing degree two by any polynomial degree does not remove
the neutral resonance obstruction. An algorithm retaining an invertible
Schr\"oder chart at every endpoint cannot cover a parameter set containing
such a fixed point. Quotient charts can avoid that representation
obstruction; their numerical convergence remains a separate problem.

\subsection{Removable branching, reconstruction and the degree limit}
\label{sec:degree-exponential}
Three structural features of this framework deserve emphasis. First,
the exchange of the two square-root fixed-point branches disappears
in $p$ and $\mathcal R$, and the apparent singularity of $\mathcal R$
is removable with value equal to the iterative residue. Second, the
quadratic equation in Corollary~\ref{cor:spectral-classification}
reconstructs both multiplier logarithms from one intrinsic spectral
function. Third, an exact affine normalisation connects polynomial
degree with the exponential map. We now prove this last connection
and follow it through spectral data, Fatou coordinates and orbit marks.

\begin{proposition}[An analytic bridge in reciprocal degree]
\label{prop:degree-bridge}
Let $x_d,c_d$ be as in Proposition~\ref{prop:unicritical-fold}, and set
\[
 A_d(u)=x_d\left(1+\frac{u}{d-1}\right),\qquad
 c=c_d-\frac{x_d}{d-1}\epsilon.
\]
Here $\epsilon$ is the rescaled cusp parameter
$(d-1)(c_d-c)/x_d$, rather than the physical difference $c_d-c$ used
in Proposition~\ref{prop:unicritical-fold}.
Then the exact affine conjugate is
\[
 A_d^{-1}\circ f_c\circ A_d(u)=G_{d,\epsilon}(u)
 =\frac{d-1}{d}\left[\left(1+\frac{u}{d-1}\right)^d-1\right]-\epsilon.
\]
Put $\tau=1/(d-1)$. The germ family has the analytic interpolation
\[
 G_{\tau,\epsilon}(u)
 =\frac{\exp\bigl((1+1/\tau)\Log(1+\tau u)\bigr)-1}{1+\tau}
       -\epsilon,\qquad
 G_{0,\epsilon}(u)=e^u-1-\epsilon.
\]
The expression at $\tau=0$ is removable. On every compact $u$ set,
for sufficiently small $|\tau|$, this family is jointly holomorphic
and has the expansion
\[
 G_{\tau,\epsilon}(u)=e^u-1-\epsilon+\tau B(u)+O(\tau^2),\qquad
 B(u)=1+e^u\left(u-1-\frac{u^2}{2}\right).
\]
The remainder and every fixed number of $u$ derivatives are uniform
on compact sets. In particular $G_{d,\epsilon}\to e^u-1-\epsilon$
at rate $O(1/(d-1))$ locally uniformly, also after differentiating.
At the cusp,
\[
 G_{\tau,0}(u)=u+\frac{u^2}{2}
 +\sum_{k\ge3}\frac{\prod_{j=1}^{k-2}(1-j\tau)}{k!}u^k,
 \qquad \rho_\tau=\frac{1+2\tau}{3}.
\]
\end{proposition}
\begin{proof}
Use $x_d^d=x_d/d$ in the affine conjugacy. For $|\tau u|<1/2$,
\[
 (1+1/\tau)\Log(1+\tau u)
 =u+\tau(u-u^2/2)+O(\tau^2),
\]
with a normally convergent power series on smaller polydiscs.
Expanding its exponential and $(1+\tau)^{-1}$ proves both
removability and the formula for $B$. Cauchy's estimates give the
derivative bounds on a slightly smaller compact set. The binomial
coefficients simplify to the displayed products. At integer degree
the product vanishes for $k>d$, as it should. The cubic coefficient
is $(1-\tau)/6$, so $1-a_3/(1/2)^2=(1+2\tau)/3$.
\end{proof}
Near $u=0$ this interpolation is holomorphic for $\tau$ in a
neighbourhood of the whole real interval $[0,1]$: choose $u$ small
enough that $|\tau u|<1/2$ and keep $\tau$ away from $-1$.
For generic $\tau$ it is only a chosen holomorphic branch, not an
entire polynomial. Global properties of the integer-degree maps
are not inherited by every interpolating parameter.

\begin{proposition}[The exponential limit is the exponential-base family]
\label{prop:degree-tetration}
For $\epsilon$ near zero set
\[
 \alpha=e^{-1-\epsilon},\qquad b=e^\alpha,\qquad
 C_\epsilon(w)=\alpha w-1-\epsilon.
\]
With the local logarithm $\log b=\alpha$, the exponential-base map
$E_b(w)=b^w$ satisfies the exact conjugacy
\[
 C_\epsilon\circ E_b\circ C_\epsilon^{-1}(u)
   =e^u-1-\epsilon=G_{0,\epsilon}(u).
\]
The cusp $\epsilon=0$ corresponds to $b=e^{1/e}$, and
$b'(0)=-e^{1/e}/e\ne0$, so this is a local analytic parameter change.
The limiting singular-value mark $-1-\epsilon$ is $C_\epsilon(0)$.
Following it twice gives $C_\epsilon(b)$, whereas the standard
exponential mark $1$ is $C_\epsilon(1)$ after conjugacy. The two
marks differ by one forward dynamical step, not by an arbitrary
replacement of the height normalisation.
\end{proposition}
\begin{proof}
Since $\alpha e^{1+\epsilon}=1$, substitution of
$w=(u+1+\epsilon)/\alpha$ gives the conjugacy. Differentiating
$b=\exp(e^{-1-\epsilon})$ gives its derivative. Finally
$E_b(0)=1$ and $E_b(1)=b$, proving all three mark identities.
\end{proof}
Thus the degree limit connects to the same exponential-base family
as the original manuscript. The limit of the parabolic marked germs
below uses the explicitly stated orbit mark; comparison with the
standard mark is by the indicated regular dynamical chart.

\begin{theorem}[The spectral bridge and its first deformation term]
\label{thm:degree-spectral}
There are a neighbourhood $V$ of $[0,1]$ and $p_0>0$ such that the
intrinsic spectral function of $G_{\tau,\epsilon}$ is jointly
holomorphic in $(\tau,p)\in V\times\{|p|<p_0\}$. It satisfies
\[
 \mathcal R_\tau(p)=\frac{1+2\tau}{3}
 +\frac{(2\tau+1)(8\tau^2+8\tau-1)}{540}\,p+O(p^2).
\]
On every smaller closed $p$ disc,
$\mathcal R_{1/(d-1)}-\mathcal R_0=O(1/(d-1))$, with the same
rate for every fixed derivative in $p$. In particular
\[
 \mathcal R_0(p)=\frac13-\frac{p}{540}+O(p^2),\qquad
 \mathcal R_{1/(d-1)}'(0)
 =\frac{(d+1)(-d^2+10d-1)}{540(d-1)^3}.
\]
The last coefficient is positive for $2\le d\le9$ and negative
for every integer $d\ge10$.
\end{theorem}
\begin{proof}
Apply the parameter-uniform fixed-point construction of
Proposition~\ref{prop:spectral-deformation} to this analytic family.
Here $a_2=1/2$ and $\gamma=1$ for every $\tau$, so
$p=2\epsilon+O(\epsilon^2)$ uniformly near the compact interval.
The inverse parameter theorem and cancellation of the two symmetric
logarithms give joint holomorphy, after choosing common neighbourhoods
by compactness. Holomorphy in $\tau$ gives the rate on smaller
polydiscs, and Cauchy's estimates give it for $p$ derivatives.

For the explicit coefficient set $\epsilon=t^2/2$. The two fixed
points have the expansion
\[
\begin{aligned}
 u_\pm={}&\pm t+\frac{\tau-1}{6}t^2
 \mp\frac{(\tau-1)(\tau+2)}{72}t^3\\
 &+\frac{(\tau-1)(\tau+2)(2\tau+1)}{540}t^4+O(t^5).
\end{aligned}
\]
For $\tau\ne0$ their multiplier logarithms are
$\ell_\pm=\tau^{-1}\Log(1+\tau u_\pm)$; at $\tau=0$ they are
$u_\pm$. Substitution gives
\[
\begin{aligned}
 \ell_\pm={}&\pm t-\frac{2\tau+1}{6}t^2
 \pm\frac{11\tau^2+11\tau+2}{72}t^3\\
 &-\frac{(2\tau+1)(43\tau^2+43\tau+4)}{1080}t^4+O(t^5),\\
 p={}&t^2+\frac{7\tau^2+7\tau+1}{36}t^4+O(t^6).
\end{aligned}
\]
The symmetric sums are even in $t$. Taking
$-(\ell_-+\ell_+)/p$ and replacing $t^2$ by $p+O(p^2)$ yields
the formula. At $\tau=1/(d-1)$ its numerator has the sign of
$-d^2+10d-1$, whose larger zero is $5+2\sqrt6\in(9,10)$.
\end{proof}
This sign change describes an intrinsic spectral coefficient. It is
not by itself a dynamical bifurcation or an algorithmic phase transition.
The companion \texttt{probe\_degree\_exponential.py} solves the full
two-fixed-point equations at fixed $p$, rather than evaluating the
displayed truncated expansion. At $p=0.01$ and 80 decimal digits it
compares degrees $2,3,9,10,20,50,100,200$ with the exponential map.
The maximum residual of its three equations is below $10^{-68}$,
and two-scale extrapolation of $\mathcal R'(0)$ differs from the
formula by less than $7\cdot10^{-16}$ in these runs. The output is
\texttt{degree-exponential-probe.json}. These figures are numerical
illustrations, not interval certificates or bounds for Fatou coordinates.
The algebraic expansion above is the proof of the spectral coefficient.
To reproduce the probe from the repository root run
\begin{verbatim}
python3 docs/probe_degree_exponential.py \
  --output docs/proofs/degree-exponential-probe.json
\end{verbatim}

\begin{theorem}[Uniform local Fatou coordinates along the bridge]
\label{thm:degree-fatou}
At $\epsilon=0$, the attracting and repelling Fatou coordinates of
$G_{\tau,0}$ can be consistently normalised on common local petals,
jointly holomorphically in $\tau$ near zero. They converge to those
of $e^u-1$ at rate $O(\tau)$ on compact subsets of these petals.
The local upper and lower horn lifts can be chosen on common end
half-planes; there they also depend holomorphically on $\tau$ and
converge at rate $O(\tau)$ on compact subsets. Every fixed Fourier
coefficient converges at this rate. If the limiting $h_1^+$ is nonzero,
so do the phase-free coordinates $J_n^\pm$ of
Proposition~\ref{prop:horn-invariants}.
\end{theorem}
\begin{proof}
Here is an explicit local coordinate construction. Set $v=-2/u$,
$g_\tau=G_{\tau,0}$ and
\[
 P_\tau(v)=-\frac2{g_\tau(-2/v)}
          =v+1+\frac{\rho_\tau}{v}+O(v^{-2}).
\]
The Laurent remainder is uniform on a closed small $\tau$ disc.
For large $R$, $P_\tau$ maps $\re v>R$ into itself and increases
the real part by at least $1/2$. Put
$a_\tau(v)=v-\rho_\tau\Log v$ there. Its Abel error
\[
 E_\tau(v)=a_\tau(P_\tau(v))-a_\tau(v)-1
\]
is uniformly $O(v^{-2})$. The normally convergent series
\[
 \Phi^{\rm att}_\tau(v)
 =a_\tau(v)+\sum_{n\ge0}E_\tau(P_\tau^{\circ n}(v))
\]
is jointly holomorphic and satisfies
$\Phi^{\rm att}_\tau\circ P_\tau=\Phi^{\rm att}_\tau+1$.
Indeed the sum is bounded by
$C\sum_{n\ge0}(\re v+n/2)^{-2}=O(1/\re v)$; subtracting its
first term proves the equation. This also fixes the asymptotic
normalisation $\Phi^{\rm att}_\tau-a_\tau\to0$ at that end.
On a farther right half-plane, Cauchy's estimate gives
$|(\Phi^{\rm att}_\tau)'-1|<1/2$, so the coordinate is univalent.

The uniform inverse germ $K_\tau=P_\tau^{-1}$ has expansion
$v-1-\rho_\tau/v+O(v^{-2})$. On $\re v<-R$, repeat the construction
with $a^-_\tau(v)=v-\rho_\tau\Log(-v)$ and
$E^-_\tau=a^-_\tau\circ K_\tau-a^-_\tau+1$.
The resulting coordinate satisfies
$\Phi^{\rm rep}_\tau\circ K_\tau=\Phi^{\rm rep}_\tau-1$ and is
univalent on a farther left half-plane. Returning by $u=-2/v$ gives
common physical petals. Holomorphy in $\tau$ proves the stated
compact-set rates by a parameter Cauchy estimate.

For the uniform end domains use the compact-family petal and horn
estimates in \cite[Propositions 29, 31 and 33]{Cheritat}, applied to
our germs restricted to a common small disc. Their quadratic
coefficient is the fixed nonzero value $1/2$, and the family on a
closed $\tau$ disc is compact in the topology of locally uniform
convergence; these are precisely the needed hypotheses.
The horn transitions are obtained from the two coordinates by local
inversion and finite holomorphic orbit charts. The uniform end
domains and local inverse theorem therefore give joint holomorphy
there; \cite[Proposition 34]{Cheritat} also supplies continuous
dependence of the extended maps on their surviving compact domains.
These are applications of existing Fatou-coordinate continuity
results, not new versions of them.
On a fixed line $t=x+iY$, $0\le x\le1$, above the common threshold,
Fourier integration of $H^+_\tau(t)-t$ proves holomorphy and the
$O(\tau)$ rate for each upper coefficient. A fixed lower line does
the same for the lower coefficients. Division by $h_1^+$ is justified
only under the nonvanishing assumption in the statement.
\end{proof}
This theorem concerns the parabolic slice $\epsilon=0$ and local horn
germs. It does not establish uniform convergence of near-parabolic
transit maps in a coupled limit $\epsilon\to0$, $d\to\infty$.

\begin{proposition}[The critical-value mark survives the degree limit]
\label{prop:degree-mark}
For integer degree the critical point and value of the normalised map
are
\[
 u_{\rm crit}=-(d-1),\qquad
 v_{d,\epsilon}=-\frac{d-1}{d}-\epsilon.
\]
The point tends to infinity, but the value tends to the asymptotic
value $-1-\epsilon$ of $e^u-1-\epsilon$.
At the cusp put
\[
 v_\tau=-\frac1{1+\tau},\quad
 m_\tau=G_{\tau,0}^{\circ2}(v_\tau),\quad
 m_0=e^{e^{-1}-1}-1.
\]
Near $\tau=0$ there is a jointly holomorphic family of parabolic
height germs $S_\tau(z)$, on a common neighbourhood of $(0,0)$, with
\[
 S_\tau(0)=m_\tau,\qquad
 \phi^{\rm att}_\tau(S_\tau(z))
   =\phi^{\rm att}_\tau(m_\tau)+z.
\]
On a smaller common height disc,
$S_{1/(d-1)}-S_0=O(1/(d-1))$, also for every fixed height derivative.
For integer degree this mark is exactly the affine image of
$f_{c_d}^3(0)$. These are singular-value marked parabolic germs;
their identification with a limit of the exterior Kneser families
is an additional question.
\end{proposition}
\begin{proof}
The derivative of $G_{d,\epsilon}$ is
$(1+u/(d-1))^{d-1}$, proving its critical point and value.
For the exponential limit, letting $u\to-\infty$ along the real
axis gives the asserted asymptotic value.
The limiting orbit from $v_0=-1$ is negative and strictly increasing
to zero: $e^u-1>u$ and $e^u-1<0$ for real $u<0$.
Some finite iterate therefore enters the common attracting petal
of Theorem~\ref{thm:degree-fatou}. Finite orbit continuity places
the nearby $v_\tau$ orbits in the same petal after the same number
of steps. All these finite limiting orbit derivatives are $e^u\ne0$;
they remain nonzero on parameter and physical neighbourhoods.
The interpolation is holomorphic on these finitely many physical
neighbourhoods for sufficiently small $|\tau|$, since
$|\tau u|<1/2$ there.
Pulling back the Fatou coordinate by those finite forward iterates
gives a jointly holomorphic regular coordinate near $m_\tau$.
Its derivative is a product of nonzero orbit derivatives and the
nonzero petal derivative. The holomorphic inverse theorem gives
$S_\tau$. Parameter holomorphy and Cauchy's estimates give the rate
and derivative convergence on a smaller common disc.
Finally $A_d^{-1}(0)=-(d-1)$ and the critical value is its first
iterate; hence $m_\tau$ is its third iterate, exactly as claimed.
\end{proof}

\begin{corollary}[First response of the marked height germ]
\label{cor:degree-response}
Let $\phi_0=\phi^{\rm att}_0$ and
$V(u)=\left.\partial_\tau\phi^{\rm att}_\tau(u)\right|_{\tau=0}$
on the regular coordinate neighbourhoods just constructed. Put
$\dot m=\left.\partial_\tau m_\tau\right|_0$ and
$W(z)=\left.\partial_\tau S_\tau(z)\right|_0$. Then
\[
 W(z)=\frac{V(m_0)+\phi_0'(m_0)\dot m-V(S_0(z))}
                  {\phi_0'(S_0(z))},\qquad
 S_{1/(d-1)}(z)=S_0(z)+\frac{W(z)}{d-1}+O((d-1)^{-2}).
\]
Writing $y=e^{-1}-1$, the mark response is explicit:
\[
 \dot m=B(y)+e^y\left(1-\frac{3}{2e}\right).
\]
The formula for $W$ is unaffected by a parameter-dependent additive
normalisation of the Fatou coordinate.
\end{corollary}
\begin{proof}
Differentiate the defining Abel-coordinate identity for $S_\tau$;
its nonzero physical derivative gives the displayed quotient.
The analytic Taylor expansion in $\tau$ gives the second formula.
Since $\dot v_0=1$, the derivative of the first marked iterate is
$B(-1)+e^{-1}=1-3/(2e)$. Differentiating the second iterate gives
$\dot m$. Adding a function of $\tau$ to $\phi_\tau$ adds the same
constant to $V(m_0)$ and $V(S_0(z))$, which cancels.
\end{proof}
Thus the map limit extends to an intrinsic spectral limit, local
horn moduli, and a regular singular-value marked parabolic height
germ with a first-order degree response. Global marked cylinders,
parameter-wide Kneser continuation and coupled near-parabolic limits
still require their own geometric and uniform orbit estimates.

\subsection{Classification and problems left by the theory}
The results above support the following precise research problems.
\begin{enumerate}
\item \emph{Membership and maximal continuation domains.} For $w^d+c$,
  $d\ge3$, construct an admissible marked atlas near the positive cusp,
  then determine how far it extends. Identify which failures are
  geometric, which are critical inverse-branch obstructions, and which
  are only failures of a chosen chart. No whole-half-plane result for
  these degrees is asserted here.
\item \emph{Classification of unfoldings.} Fix the simple parabolic
  formal type and its horn pair. Decide which deformations are
  distinguished by the joint data $\mathcal R(p)$ and $j_{n,r}^\pm$,
  whether the latter exist in a useful class, and whether additional
  sector transition data are necessary. Proposition~\ref{prop:same-cusp-different-unfolding}
  already distinguishes deformations with identical cusp germs.
  The pointwise classification in
  Proposition~\ref{prop:horn-invariants} does not answer this question.
\item \emph{Several singular orbits.} Establish cylinder-entry
  hypotheses for a family with several independently tracked singular
  orbits. Determine which relative-mark functions and braid representations
  are realisable and whether they contribute to global classification.
\item \emph{Deformation stability.} Prove openness of admissibility in
  a specified topology with specified domains.
  Proposition~\ref{prop:finite-witness-stability} already gives explicit
  stability inequalities for finite witnesses on a compact core,
  but supplies neither the two infinite ends nor a noncompact tail.
  In particular global lower-half-plane injectivity of a quadratic
  must not be assumed stable under a higher-degree perturbation.
\item \emph{Certified neutral numerical charts.} Approximate the marked
  quotient itself, compare it on overlaps with nonneutral end charts,
  and prove an error bound including uniformisation, inverse branches,
  truncation and arithmetic. Establishing such a method would add a
  convergence result that Theorem~\ref{thm:general-resonance} excludes
  for an everywhere-Schr\"oder representation.
\item \emph{Degree and cusp limits together.} Theorems
  \ref{thm:degree-spectral} and~\ref{thm:degree-fatou} prove the local
  spectral and parabolic-coordinate limits; Proposition~\ref{prop:degree-mark}
  preserves the singular-value orbit mark. Determine uniform
  near-parabolic transit estimates when $\tau\to0$ and $\epsilon\to0$
  together, and compare the resulting sewn or exterior Kneser germs.
  The finite-index orbit estimates used above do not control transit
  times that diverge as the cusp is approached.
\end{enumerate}
The abstract continuation theorem, coordinate identities, spectral
deformation invariant and separating example, polynomial fold
calculation and resonance theorem are proved statements here.
The reciprocal-degree interpolation, spectral expansion and marked
parabolic germ limit are also proved; the horn convergence applies
the stated existing compact-family Fatou-coordinate results.
The global membership, realisation, deformation completeness and
numerical convergence questions remain open in this framework.

\begin{thebibliography}{9}
\bibitem{AB} L. Ahlfors and L. Bers, \emph{Riemann's mapping theorem for variable
metrics}, Ann. of Math. 72 (1960), 385--404.
\url{https://www.cimat.mx/~mmoreno/teaching/fall10/AhlforsBers.pdf}.
\bibitem{TK} H. Trappmann and D. Kouznetsov, \emph{Uniqueness of holomorphic
Abel functions at a complex fixed point pair}, Aequationes Math. 81 (2011).
\url{https://arxiv.org/abs/1006.3981}.
\bibitem{BMS} A. Berretti, S. Marmi and D. Sauzin,
\emph{Limit at resonances of linearizations of some complex analytic
dynamical systems}, Ergodic Theory Dynam. Systems 20 (2000), 963--990.
\url{https://homepage.sns.it/marmi/papers/bms.pdf}.
\bibitem{Cheritat} A. Ch\'eritat, \emph{Near parabolic renormalization for
unicritical holomorphic maps}, Arnold Math. J. 8 (2022), 169--270.
\url{https://arxiv.org/abs/1404.4735}.
\bibitem{DS} A. Dudko and D. Sauzin, \emph{The resurgent character of the
Fatou coordinates of a simple parabolic germ}, C. R. Math. 352 (2014),
255--261. \url{https://arxiv.org/abs/1307.8093}.
\bibitem{local} The repository manuscript, \texttt{docs/paper-submission/main.tex},
Section ``Continuation of Kneser's solution around the cusp''.
The quadratic orbit, inverse charts, transport field, overlap comparison,
and sewn branch are established above; exponential-family continuation
theorems are not assumed.
\end{thebibliography}
\end{document}
